\documentclass[aspectratio=169]{beamer}\usepackage[]{graphicx}\usepackage[]{xcolor}
% maxwidth is the original width if it is less than linewidth
% otherwise use linewidth (to make sure the graphics do not exceed the margin)
\makeatletter
\def\maxwidth{ %
  \ifdim\Gin@nat@width>\linewidth
    \linewidth
  \else
    \Gin@nat@width
  \fi
}
\makeatother

\definecolor{fgcolor}{rgb}{0.345, 0.345, 0.345}
\newcommand{\hlnum}[1]{\textcolor[rgb]{0.686,0.059,0.569}{#1}}%
\newcommand{\hlsng}[1]{\textcolor[rgb]{0.192,0.494,0.8}{#1}}%
\newcommand{\hlcom}[1]{\textcolor[rgb]{0.678,0.584,0.686}{\textit{#1}}}%
\newcommand{\hlopt}[1]{\textcolor[rgb]{0,0,0}{#1}}%
\newcommand{\hldef}[1]{\textcolor[rgb]{0.345,0.345,0.345}{#1}}%
\newcommand{\hlkwa}[1]{\textcolor[rgb]{0.161,0.373,0.58}{\textbf{#1}}}%
\newcommand{\hlkwb}[1]{\textcolor[rgb]{0.69,0.353,0.396}{#1}}%
\newcommand{\hlkwc}[1]{\textcolor[rgb]{0.333,0.667,0.333}{#1}}%
\newcommand{\hlkwd}[1]{\textcolor[rgb]{0.737,0.353,0.396}{\textbf{#1}}}%
\let\hlipl\hlkwb

\usepackage{framed}
\makeatletter
\newenvironment{kframe}{%
 \def\at@end@of@kframe{}%
 \ifinner\ifhmode%
  \def\at@end@of@kframe{\end{minipage}}%
  \begin{minipage}{\columnwidth}%
 \fi\fi%
 \def\FrameCommand##1{\hskip\@totalleftmargin \hskip-\fboxsep
 \colorbox{shadecolor}{##1}\hskip-\fboxsep
     % There is no \\@totalrightmargin, so:
     \hskip-\linewidth \hskip-\@totalleftmargin \hskip\columnwidth}%
 \MakeFramed {\advance\hsize-\width
   \@totalleftmargin\z@ \linewidth\hsize
   \@setminipage}}%
 {\par\unskip\endMakeFramed%
 \at@end@of@kframe}
\makeatother

\definecolor{shadecolor}{rgb}{.97, .97, .97}
\definecolor{messagecolor}{rgb}{0, 0, 0}
\definecolor{warningcolor}{rgb}{1, 0, 1}
\definecolor{errorcolor}{rgb}{1, 0, 0}
\newenvironment{knitrout}{}{} % an empty environment to be redefined in TeX

\usepackage{alltt}

% Set lecture number for later use


% Part common to all the lectures
\subtitle{MATH 2740 -- Mathematics of Data Science -- Unit U08}
\author{\texorpdfstring{Julien Arino\newline\url{julien.arino@umanitoba.ca}}{Julien Arino}}
\institute{Department of Mathematics @ University of Manitoba}

% Title of the lecture
\title{Non-negative matrices, digraphs and regular Markov chains}



\usetheme{default}
% Slide setup, colour independent

\usepackage{amsmath,amssymb,amsthm}
\usepackage[utf8]{inputenc}
\usepackage{colortbl}
\usepackage{bm}
\usepackage{xcolor}
\usepackage{dsfont}
\usepackage{setspace}
% To use \ding{234} and the like
\usepackage{pifont}
% To cross reference between slide files
\usepackage{zref-xr,zref-user}
% Use something like
% \zexternaldocument{fileI}
% in the tex files. And cite using \zref instead of \ref

% Cross-reference system - see CROSS-REFERENCE-SETUP.md for manual setup instructions
\usepackage{booktabs}
\usepackage{marvosym}
\usepackage{cancel}
%\usepackage{transparent}
% Make doi clickable in the bibliography?
\usepackage{doi}

\usepackage[T1]{fontenc}

\usepackage{longtable}

% For heavier titles
\usepackage{helvet} % Enables Helvetica font family


% Fields and the like
\def\IC{\mathbb{C}}
\def\IE{\mathbb{E}}
\def\IF{\mathbb{F}}
\def\II{\mathbb{I}}
\def\IJ{\mathbb{J}}
\def\IK{\mathbb{K}}
\def\IM{\mathbb{M}}
\def\IN{\mathbb{N}}
\def\IP{\mathbb{P}}
\def\IR{\mathbb{R}}
\newcommand{\IRplus}{\mathbb{R}_{\ge 0}}
\def\IZ{\mathbb{Z}}
\def\11{\mathds{1}}


% Bold lowercase
\def\ba{\bm{a}}
\def\bb{\bm{b}}
\def\bc{\bm{c}}
\def\bd{\bm{d}}
\def\be{\bm{e}}
\def\bf{\bm{f}}
\def\bg{\bm{g}}
\def\bh{\bm{h}}
\def\bi{\bm{i}}
\def\bj{\bm{j}}
\def\bk{\bm{k}}
\def\bn{\bm{n}}
\def\bp{\bm{p}}
\def\br{\bm{r}}
\def\bs{\bm{s}}
\def\bu{\bm{u}}
\def\bv{\bm{v}}
\def\bw{\bm{w}}
\def\bx{\bm{x}}
\def\by{\bm{y}}
\def\bz{\bm{z}}
\newcommand{\vect}[1]{\bm{#1}}

% Bold capitals
\def\bB{\bm{B}}
\def\bD{\bm{D}}
\def\bE{\bm{E}}
\def\bF{\bm{F}}
\def\bG{\bm{G}}
\def\bI{\bm{I}}
\def\bL{\bm{L}}
\def\bN{\bm{N}}
\def\bP{\bm{P}}
\def\bR{\bm{R}}
\def\bS{\bm{S}}
\def\bT{\bm{T}}
\def\bX{\bm{X}}

% Bold numbers
\def\b0{\bm{0}}

% Bold greek
\bmdefine{\bmu}{\bm{\mu}}
\def\bphi{\bm{\phi}}
\def\bvarphi{\bm{\varphi}}
\def\bPi{\bm{\Pi}}
\def\bGamma{\bm{\Gamma}}

% Bold red sentence
\def\boldred#1{{\color{red}\textbf{#1}}}
\def\defword#1{{\color{orange}\textbf{#1}}}

% Caligraphic letters
\def\A{\mathcal{A}}
\def\B{\mathcal{B}}
\def\C{\mathcal{C}}
\def\D{\mathcal{D}}
\def\E{\mathcal{E}}
\def\F{\mathcal{F}}
\def\G{\mathcal{G}}
\def\H{\mathcal{H}}
\def\I{\mathcal{I}}
\def\L{\mathcal{L}}
\def\M{\mathcal{M}}
\def\N{\mathcal{N}}
\def\P{\mathcal{P}}
\def\R{\mathcal{R}}
\def\S{\mathcal{S}}
\def\T{\mathcal{T}}
\def\U{\mathcal{U}}
\def\V{\mathcal{V}}

% Adding space for prime (') where needed
\def\pprime{\,'}
% Adding space for star (\star) where needed
\def\pstar{{\,\star}}

% tt font for code
\def\code#1{{\tt #1}}

% i.e., e.g.
\def\eg{\emph{e.g.}}
\def\ie{\emph{i.e.}}


% Operators and special symbols
\def\nbOne{{\mathchoice {\rm 1\mskip-4mu l} {\rm 1\mskip-4mu l}
{\rm 1\mskip-4.5mu l} {\rm 1\mskip-5mu l}}}
\def\cov{\ensuremath{\mathsf{cov}}}
\def\Var{\ensuremath{\mathsf{Var}\ }}
\def\Im{\textrm{Im}\;}
\def\Re{\textrm{Re}\;}
\def\det{\ensuremath{\mathsf{det}}}
\def\diag{\ensuremath{\mathsf{diag}}}
\def\nullspace{\ensuremath{\mathsf{null}}}
\def\nullity{\ensuremath{\mathsf{nullity}}}
\def\rank{\ensuremath{\mathsf{rank}}}
\def\range{\ensuremath{\mathsf{range}}}
\def\sgn{\ensuremath{\mathsf{sgn}}}
\def\Span{\ensuremath{\mathsf{span}}}
\def\tr{\ensuremath{\mathsf{tr}}}
\def\imply{$\Rightarrow$}
\def\restrictTo#1#2{\left.#1\right|_{#2}}
\newcommand{\parallelsum}{\mathbin{\!/\mkern-5mu/\!}}
\def\dsum{\mathop{\displaystyle \sum }}%
\def\dind#1#2{_{\substack{#1\\ #2}}}

\newcommand{\Qmatrix}[1]{%
  \begin{pmatrix}#1\end{pmatrix}%
}

\DeclareMathOperator{\GL}{GL}
\DeclareMathOperator{\Rel}{Re}
\def\Nt#1{\left|\!\left|\!\left|#1\right|\!\right|\!\right|}
\newcommand{\tripbar}{|\! |\! |}



% The beamer bullet (in base colour)
\def\bbullet{\leavevmode\usebeamertemplate{itemize item}\ }

% Theorems and the like
\newtheorem{importantdefinition}[theorem]{Definition}
\newtheorem{proposition}[theorem]{Proposition}
\newtheorem{property}[theorem]{Property}
\newtheorem{importantproperty}[theorem]{Property}
\newtheorem{importanttheorem}[theorem]{Theorem}
%\newtheorem{lemma}[theorem]{Lemma}
%\newtheorem{corollary}[theorem]{Corollary}
\newtheorem{remark}[theorem]{Remark}
\setbeamertemplate{theorems}[numbered]
%\setbeamertemplate{theorems}[ams style]

%
%\usecolortheme{orchid}
%\usecolortheme{orchid}

\def\red{\color[rgb]{1,0,0}}
\def\blue{\color[rgb]{0,0,1}}
\def\green{\color[rgb]{0,1,0}}

% Fix skipping lines after items in the bibliography
\setbeamertemplate{bibliography entry title}{}
\setbeamertemplate{bibliography entry location}{}
\setbeamertemplate{bibliography entry note}{}

% Get rid of navigation stuff
\setbeamertemplate{navigation symbols}{}

% Set footline/header line
\setbeamertemplate{footline}
{%
\quad p. \insertpagenumber \quad--\quad \insertsection\vskip2pt
}
% \setbeamertemplate{headline}
% {%
% \quad\insertsection\hfill p. \insertpagenumber\quad\mbox{}\vskip2pt
% }


\makeatletter
\newlength\beamerleftmargin
\setlength\beamerleftmargin{\Gm@lmargin}
\makeatother

% Colours for special pages
\def\extraContent{yellow!20}


%%%%%%%%%%%%%%%%%
\usepackage{tikz}
\usetikzlibrary{shapes,arrows}
\usetikzlibrary{positioning}
\usetikzlibrary{shapes.symbols,shapes.callouts,patterns}
\usetikzlibrary{calc,fit}
\usetikzlibrary{backgrounds}
\usetikzlibrary{decorations.pathmorphing,fit,petri}
\usetikzlibrary{automata}
\usetikzlibrary{fadings}
\usetikzlibrary{patterns,hobby}
\usetikzlibrary{backgrounds,fit,petri}
\usetikzlibrary{tikzmark}

\usepackage{pgfplots}
\pgfplotsset{compat=1.6}
\pgfplotsset{ticks=none}

\usetikzlibrary{decorations.markings}
\usetikzlibrary{arrows.meta}
\tikzset{>=stealth}

% For tikz
\tikzstyle{cloud} = [draw, ellipse,fill=red!20, node distance=0.87cm,
minimum height=2em]
\tikzstyle{line} = [draw, -latex']


%%% For max frame images
\newenvironment{changemargin}[2]{%
\begin{list}{}{%
\setlength{\topsep}{0pt}%
\setlength{\leftmargin}{#1}%
\setlength{\rightmargin}{#2}%
\setlength{\listparindent}{\parindent}%
\setlength{\itemindent}{\parindent}%
\setlength{\parsep}{\parskip}%
}%
\item[]}{\end{list}}


% Make one image take up the entire slide content area in beamer,.:
% centered/centred full-screen image, with title:
% This uses the whole screen except for the 1cm border around it
% all. 128x96mm
\newcommand{\titledFrameImage}[2]{
\begin{frame}{#1}
%\begin{changemargin}{-1cm}{-1cm}
\begin{center}
\includegraphics[width=108mm,height=\textheight,keepaspectratio]{#2}
\end{center}
%\end{changemargin}
\end{frame}
}

% Make one image take up the entire slide content area in beamer.:
% centered/centred full-screen image, no title:
% This uses the whole screen except for the 1cm border around it
% all. 128x96mm
\newcommand{\plainFrameImage}[1]{
\begin{frame}[plain]
%\begin{changemargin}{-1cm}{-1cm}
\begin{center}
\includegraphics[width=108mm,height=76mm,keepaspectratio]{#1}
\end{center}
%\end{changemargin}
\end{frame}
}

% Make one image take up the entire slide area, including borders, in beamer.:
% centered/centred full-screen image, no title:
% This uses the entire whole screen
\newcommand{\maxFrameImage}[1]{
\begin{frame}[plain]
\begin{changemargin}{-1cm}{-1cm}
\begin{center}
\includegraphics[width=\paperwidth,height=\paperheight,keepaspectratio]
{#1}
\end{center}
\end{changemargin}
\end{frame}
}

% This uses the entire whole screen (to include in frame)
\newcommand{\maxFrameImageNoFrame}[1]{
\begin{changemargin}{-1cm}{-1cm}
\begin{center}
\includegraphics[width=\paperwidth,height=0.99\paperheight,keepaspectratio]
{#1}
\end{center}
\end{changemargin}
}

% Make one image take up the entire slide area, including borders, in beamer.:
% centered/centred full-screen image, no title:
% This uses the entire whole screen
\newcommand{\maxFrameImageColor}[2]{
\begin{frame}[plain]
\setbeamercolor{normal text}{bg=#2!20}
\begin{changemargin}{-1cm}{-1cm}
\begin{center}
\includegraphics[width=\paperwidth,height=\paperheight,keepaspectratio]
{#1}
\end{center}
\end{changemargin}
\end{frame}
}


\usepackage{tikz}
\usetikzlibrary{patterns,hobby,matrix}
\usepackage{pgfplots}
\pgfplotsset{compat=1.6}
\pgfplotsset{ticks=none}

\usetikzlibrary{backgrounds}
\usetikzlibrary{decorations.markings}
\usetikzlibrary{arrows.meta}
\tikzset{>=stealth}

\tikzset{
  clockwise arrows/.style={
    postaction={
      decorate,
      decoration={
        markings,
        mark=between positions 0.1 and 0.9 step 40pt with {\arrow{>}},
   }}}}


% New integrated section command: creates section and section slide
\newcommand{\Ssection}[2]{
\section{#1}
\begin{frame}[noframenumbering,plain]
  \begin{tikzpicture}[remember picture,overlay]
    \node[above right,inner sep=0pt,opacity=0.2] at (current page.south west)
    {
        \includegraphics[height=\paperheight,width=\paperwidth]{#2}
    };
  \end{tikzpicture}
  \setbeamercolor{section in toc}{fg=section_page_list_colour}
  \setbeamerfont{section in toc}{size=\Large,series=\bfseries}
  \setbeamertemplate{section in toc shaded}[default][60]
  \tableofcontents[
    currentsection,
    sectionstyle=show/shaded,
    subsectionstyle=show/hide/hide,
    subsubsectionstyle=hide/hide/hide]
\end{frame}
\addtocounter{page}{-1}
}

% New integrated section command with subsections: creates section and section slide showing subsections
\newcommand{\SsectionWithSubs}[2]{
\section{#1}
\begin{frame}[noframenumbering,plain]
  \begin{tikzpicture}[remember picture,overlay]
    \node[above right,inner sep=0pt,opacity=0.2] at (current page.south west)
    {
        \includegraphics[height=\paperheight,width=\paperwidth]{#2}
    };
  \end{tikzpicture}
  \setbeamercolor{section in toc}{fg=section_page_list_colour}
  \setbeamerfont{section in toc}{size=\Large,series=\bfseries}
  \setbeamertemplate{section in toc shaded}[default][60]
  \tableofcontents[
    currentsection,
    sectionstyle=show/hide,
    subsectionstyle=show/show/hide,
    subsubsectionstyle=hide/hide/hide]
\end{frame}
\addtocounter{page}{-1}
}

% New integrated subsection command: creates subsection and subsection slide
\newcommand{\Ssubsection}[2]{
\subsection{#1}
\begin{frame}[noframenumbering,plain]
  \begin{tikzpicture}[remember picture,overlay]
    \node[above right,inner sep=0pt,opacity=0.2] at (current page.south west)
    {
        \includegraphics[height=\paperheight,width=\paperwidth]{#2}
    };
  \end{tikzpicture}
  \setbeamercolor{section in toc}{fg=subsection_page_list_colour}
  \setbeamerfont{section in toc}{size=\Large,series=\bfseries}
  \setbeamertemplate{section in toc shaded}[default][60]
  \setbeamerfont{subsection in toc}{series=\bfseries}
  \setbeamertemplate{subsection in toc shaded}[default][50]
  \tableofcontents[
    currentsection,
    sectionstyle=show/hide,
    subsectionstyle=show/shaded/hide,
    subsubsectionstyle=hide/hide/hide]
\end{frame}
\addtocounter{page}{-1}
}

% New integrated subsubsection command: creates subsubsection and subsubsection slide
\newcommand{\Ssubsubsection}[2]{
\subsubsection{#1}
\begin{frame}[noframenumbering,plain]
  \begin{tikzpicture}[remember picture,overlay]
    \node[above right,inner sep=0pt,opacity=0.2] at (current page.south west)
    {
        \includegraphics[height=\paperheight,width=\paperwidth]{#2}
    };
  \end{tikzpicture}
  \setbeamercolor{section in toc}{fg=subsub_header_section}
  \setbeamerfont{section in toc}{size=\Large,series=\bfseries}
  \setbeamertemplate{section in toc shaded}[default][60]
  \setbeamerfont{subsection in toc}{series=\bfseries}
  \setbeamertemplate{subsection in toc shaded}[default][50]
  \setbeamertemplate{subsubsection in toc shaded}[default][50]
  \tableofcontents[
    currentsection,
    sectionstyle=show/hide,
    subsectionstyle=show/hide/hide,
    subsubsectionstyle=show/shaded/hide]
\end{frame}
\addtocounter{page}{-1}
}

% Legacy commands (kept for backward compatibility)
% Beginning of a section
\newcommand{\newSectionSlide}[1]{
\begin{frame}[noframenumbering,plain]
  \begin{tikzpicture}[remember picture,overlay]
    \node[above right,inner sep=0pt,opacity=0.2] at (current page.south west)
    {
        \includegraphics[height=\paperheight,width=\paperwidth]{#1}
    };
  \end{tikzpicture}
  \setbeamercolor{section in toc}{fg=section_page_list_colour}
  \setbeamerfont{section in toc}{size=\Large,series=\bfseries}
  \setbeamertemplate{section in toc shaded}[default][60]
  \tableofcontents[
    currentsection,
    sectionstyle=show/shaded,
    subsectionstyle=show/hide/hide,
    subsubsectionstyle=hide/hide/hide]
\end{frame}
\addtocounter{page}{-1}
}

% Beginning of a section in which we also show subsections
\newcommand{\newSectionWithSubsSlide}[1]{
	\begin{frame}[noframenumbering,plain]
		\begin{tikzpicture}[remember picture,overlay]
			\node[above right,inner sep=0pt,opacity=0.2] at (current page.south west)
			{
				\includegraphics[height=\paperheight,width=\paperwidth]{#1}
			};
		\end{tikzpicture}
		\setbeamercolor{section in toc}{fg=section_page_list_colour}
		\setbeamerfont{section in toc}{size=\Large,series=\bfseries}
		\setbeamertemplate{section in toc shaded}[default][60]
		\tableofcontents[
		currentsection,
		sectionstyle=show/hide,
		subsectionstyle=show/show/hide,
		subsubsectionstyle=hide/hide/hide]
	\end{frame}
	\addtocounter{page}{-1}
}

% Beginning of a subsection
\newcommand{\newSubSectionSlide}[1]{
\begin{frame}[noframenumbering,plain]
  \begin{tikzpicture}[remember picture,overlay]
    \node[above right,inner sep=0pt,opacity=0.2] at (current page.south west)
    {
        \includegraphics[height=\paperheight,width=\paperwidth]{#1}
    };
  \end{tikzpicture}
  \setbeamercolor{section in toc}{fg=subsection_page_list_colour}
  \setbeamerfont{section in toc}{size=\Large,series=\bfseries}
  \setbeamertemplate{section in toc shaded}[default][60]
  \setbeamerfont{subsection in toc}{series=\bfseries}
  \setbeamertemplate{subsection in toc shaded}[default][50]
  \tableofcontents[
    currentsection,
    sectionstyle=show/hide,
    subsectionstyle=show/shaded/hide,
    subsubsectionstyle=hide/hide/hide]
\end{frame}
\addtocounter{page}{-1}
}

% Beginning of a subsubsection
\newcommand{\newSubSubSectionSlide}[1]{
\begin{frame}[noframenumbering,plain]
  \begin{tikzpicture}[remember picture,overlay]
    \node[above right,inner sep=0pt,opacity=0.2] at (current page.south west)
    {
        \includegraphics[height=\paperheight,width=\paperwidth]{#1}
    };
  \end{tikzpicture}
  \setbeamercolor{section in toc}{fg=subsub_header_section}
  \setbeamerfont{section in toc}{size=\Large,series=\bfseries}
  \setbeamertemplate{section in toc shaded}[default][60]
  \setbeamerfont{subsection in toc}{series=\bfseries}
  \setbeamertemplate{subsection in toc shaded}[default][50]
  \setbeamertemplate{subsubsection in toc shaded}[default][50]
  \tableofcontents[
    currentsection,
    sectionstyle=show/hide,
    subsectionstyle=show/hide/hide,
    subsubsectionstyle=show/shaded/hide]
\end{frame}
\addtocounter{page}{-1}
}


   %%%%%%%%%%%
% To have links to parts in the outline
\makeatletter
\AtBeginPart{%
  \addtocontents{toc}{\protect\beamer@partintoc{\the\c@part}{\beamer@partnameshort}{\the\c@page}}%
}
%% number, shortname, page.
\providecommand\beamer@partintoc[3]{%
  \ifnum\c@tocdepth=-1\relax
    % requesting onlyparts.
    \makebox[6em]{Part #1:} \textcolor{green!30!blue}{\hyperlink{#2}{#2}}
    \par
  \fi
}
\define@key{beamertoc}{onlyparts}[]{%
  \c@tocdepth=-1\relax
}
\makeatother%

\newcommand{\nameofthepart}{}
\newcommand{\nupart}[1]%
    {   \part{#1}%
        \renewcommand{\nameofthepart}{#1}%
        {
          \setbeamercolor{background canvas}{bg=orange!50}
          \begin{frame}{#1}%\partpage 
          \hypertarget{\nameofthepart}{}\tableofcontents%
          \end{frame}
        }
    }

% This command creates a title page using TikZ only
\newcommand{\tikztitlepage}[1]{%
\begin{frame}[plain,noframenumbering]
  \begin{tikzpicture}[remember picture,overlay]
    % Background image
    \node[above right,inner sep=0pt,opacity=0.1] 
      at (current page.south west) 
      {\includegraphics[width=\paperwidth,height=\paperheight]{#1}};

    % University logo
    \node[anchor=north east, inner sep=5pt, opacity=0.9] 
      at (current page.north east)
      {\includegraphics[width=0.2\textwidth]{FIGS-slides-admin/UM-logo-horizontal-CMYK.png}};
    
    % Title
    \node[anchor=center, align=center, 
          font=\fontsize{13}{15}\bfseries\color{UMbrown}, 
          text width=0.9\textwidth] 
          at ([yshift=2cm]current page.center)
          {\inserttitle};

      % Authors
      \node[anchor=center, align=center,
        font=\fontsize{10}{12}\bfseries\color{UMbrown},
        text width=0.7\textwidth]
        at ([yshift=0.8cm]current page.center)
        {\insertauthor};

      % Affiliation
      \node[anchor=north, align=center,
        font=\fontsize{9}{11}\color{UMbrown},
        text width=0.7\textwidth]
        at ([yshift=-0.2cm]current page.center)
        {\insertaffiliation};      
    % Date
    \node[anchor=north, align=center, 
          font=\fontsize{12}{16}\bfseries\color{UMbrown},
          text width=0.7\textwidth] 
          at ([yshift=0.2cm]current page.center)
          {\insertdate};

    % Land acknowledgement
    \node[anchor=south, align=justify, 
          font=\footnotesize, text=black, 
          text width=1.1\textwidth] 
          at ([yshift=0.5cm]current page.south)
          {The University of Manitoba campuses are located on original lands of Anishinaabeg, Ininew, Anisininew, Dakota and Dene peoples, and on the National Homeland of the Red River Métis.\\
          We respect the Treaties that were made on these territories, we acknowledge the harms and mistakes of the past, and we dedicate ourselves to move forward in partnership with Indigenous communities in a spirit of Reconciliation and collaboration.};
  \end{tikzpicture}
  \addtocounter{page}{-1}
\end{frame}
}
% The title page with figure
% \newcommand{\titlepagewithfigure}[1]{%
%   \begin{frame}[noframenumbering,plain]
%     \begin{tikzpicture}[remember picture,overlay]
%       \node[above right,inner sep=0pt,opacity=0.1] at (current page.south west)
%       {
%           \includegraphics[height=\paperheight,width=\paperwidth]{#1}
%       };
%       \node[anchor=north east,
%       inner sep=5pt,
%       opacity=0.9] at (current page.north east)
%       {
%           \includegraphics[width=0.2\textwidth]{FIGS-slides-admin/UM-logo-horizontal-CMYK.png}
%       };
%       \node[anchor=south, 
%       align=justify, 
%       text=black, 
%       text width=1.1\textwidth,
%       font=\footnotesize]  (land_acknowledgement)
%       at (current page.south) 
%       {The University of Manitoba campuses are located on original lands of Anishinaabeg, Ininew, Anisininew, Dakota and Dene peoples, and on the National Homeland of the Red River Métis.
%       We respect the Treaties that were made on these territories, we acknowledge the harms and mistakes of the past, and we dedicate ourselves to move forward in partnership with Indigenous communities in a spirit of Reconciliation and collaboration.};  
%       % \node[align=center, anchor=south,
%       % above=0.5cm of land_acknowledgement,
%       % text=black,
%       % font=\bfseries] {\insertdate};
%   \end{tikzpicture}
%   \setbeamercolor{title}{fg=title_page_title_colour}
%   \setbeamerfont{title}{size=\Large,series=\bfseries}
%   \setbeamercolor{author}{fg=title_page_author_colour}
%   \setbeamerfont{author}{size=\large,series=\bfseries}
%   \setbeamercolor{institute}{fg=title_page_institute_colour}
%   \setbeamerfont{institute}{size=\large,series=\bfseries}
%   \setbeamercolor{date}{fg=title_page_date_colour}
%   \setbeamerfont{date}{series=\bfseries}
% 	\titlepage
% \end{frame}
% \addtocounter{page}{-1}
% }

\newcommand{\titlepagewithfigure}[1]{%
  \begin{frame}[noframenumbering,plain]
    \begin{tikzpicture}[remember picture,overlay]
      \node[above right,inner sep=0pt,opacity=0.1] at (current page.south west)
      {
          \includegraphics[height=\paperheight,width=\paperwidth]{#1}
      };
      \node[anchor=north east,
      inner sep=5pt,
      opacity=0.9] at (current page.north east)
      {
          \includegraphics[width=0.2\textwidth]{FIGS-slides-admin/UM-logo-horizontal-CMYK.png}
      };
      \node[anchor=south, 
      align=justify, 
      text=black, 
      text width=1.1\textwidth,
      font=\footnotesize]  (land_acknowledgement)
      at (current page.south) 
      {The University of Manitoba campuses are located on original lands of Anishinaabeg, Ininew, Anisininew, Dakota and Dene peoples, and on the National Homeland of the Red River Métis.
      We respect the Treaties that were made on these territories, we acknowledge the harms and mistakes of the past, and we dedicate ourselves to move forward in partnership with Indigenous communities in a spirit of Reconciliation and collaboration.};  
      % \node[align=center, anchor=south,
      % above=0.5cm of land_acknowledgement,
      % text=black,
      % font=\bfseries] {\insertdate};
  \end{tikzpicture}
  \setbeamercolor{title}{fg=title_page_title_colour}
  \setbeamerfont{title}{size=\Large,series=\bfseries,family=\usefont{T1}{phv}{b}{n}}
  \setbeamercolor{author}{fg=title_page_author_colour}
  \setbeamerfont{author}{size=\large,series=\bfseries,family=\usefont{T1}{phv}{b}{n}}
  \setbeamercolor{institute}{fg=title_page_institute_colour}
  \setbeamerfont{institute}{size=\large,series=\bfseries,family=\usefont{T1}{phv}{b}{n}}
  \setbeamercolor{date}{fg=title_page_date_colour}
  \setbeamerfont{date}{series=\bfseries,family=\usefont{T1}{phv}{b}{n}}
	\titlepage
\end{frame}
\addtocounter{page}{-1}
}
% The outline page, with figure
% \newcommand{\outlinepage}[1]{%
% \begin{frame}[noframenumbering,plain]
%   \begin{tikzpicture}[remember picture,overlay]
%     \node[above right,inner sep=0pt,opacity=0.2] at (current page.south west)
%     {
%         \includegraphics[height=\paperheight,width=\paperwidth]{#1}
%     };
%   \end{tikzpicture}
%   \setbeamercolor{section in toc}{fg=outline_page_list_colour}
%   \setbeamerfont{section in toc}{size=\Large,series=\bfseries,family=\sffamily}
%   \frametitle{\textcolor{outline_page_title_colour}{\LARGE\bfseries Outline}}
%   \tableofcontents[hideallsubsections]
% \end{frame}
% \addtocounter{page}{-1}
% }
% The outline page, with figure
\newcommand{\outlinepage}[1]{%
\begin{frame}[noframenumbering,plain]
  \begin{tikzpicture}[remember picture,overlay]
    \node[above right,inner sep=0pt,opacity=0.2] at (current page.south west)
    {
        \includegraphics[height=\paperheight,width=\paperwidth]{#1}
    };
  \end{tikzpicture}
  \setbeamercolor{section in toc}{fg=outline_page_list_colour}
  % Use Helvetica Bold only for the outline slide TOC
  \setbeamerfont{section in toc}{size=\Large,family=\usefont{T1}{phv}{b}{n}}
  % Use Helvetica Bold for the outline title
  \frametitle{\textcolor{outline_page_title_colour}{\usefont{T1}{phv}{b}{n}\LARGE Outline}}
  \tableofcontents[hideallsubsections]
\end{frame}
\addtocounter{page}{-1}
}


%\let\oldsection\section
%\renewcommand{\section}[2]{\oldsection[#1]\newSectionSlide[#2]}


%%%%%%%%%%%%%%%%%%%%%
% CUSTOM SLIDE BACKGROUNDS
%%%%%%%%%%%%%%%%%%%%%
% Define custom background templates for different colors
\defbeamertemplate*{background canvas}{blue}{%
  \color{blue!15}\vrule width\paperwidth height\paperheight%
}
\defbeamertemplate*{background canvas}{green}{%
  \color{green!15}\vrule width\paperwidth height\paperheight%
}
\defbeamertemplate*{background canvas}{red}{%
  \color{red!15}\vrule width\paperwidth height\paperheight%
}
\defbeamertemplate*{background canvas}{yellow}{%
  \color{yellow!20}\vrule width\paperwidth height\paperheight%
}
\defbeamertemplate*{background canvas}{purple}{%
  \color{purple!15}\vrule width\paperwidth height\paperheight%
}
\defbeamertemplate*{background canvas}{orange}{%
  \color{orange!20}\vrule width\paperwidth height\paperheight%
}

% Define keys for the different background options
\makeatletter
\define@key{beamerframe}{blue}[true]{\setbeamertemplate{background canvas}[blue]}
\define@key{beamerframe}{green}[true]{\setbeamertemplate{background canvas}[green]}
\define@key{beamerframe}{red}[true]{\setbeamertemplate{background canvas}[red]}
\define@key{beamerframe}{yellow}[true]{\setbeamertemplate{background canvas}[yellow]}
\define@key{beamerframe}{purple}[true]{\setbeamertemplate{background canvas}[purple]}
\define@key{beamerframe}{orange}[true]{\setbeamertemplate{background canvas}[orange]}
\makeatother

% Reset to normal background for all frames by default
\BeforeBeginEnvironment{frame}{\setbeamertemplate{background canvas}[default]}

\newcommand{\punnett}[2]{
    \begin{center}
    \renewcommand{\arraystretch}{1.5} % Add space to rows
    \begin{tabular}{|c c | c|c|}
        \multicolumn{2}{c}{} & \multicolumn{2}{c}{\textbf{Father}} \\
        \multicolumn{2}{c}{} & #1 \\ \hline
        #2
    \end{tabular}
    \end{center}
}

% Colour definitions for Punnett squares
\usepackage{multirow}
\usepackage{colortbl}
\colorlet{punnettorange}{orange!30}
\colorlet{punnettblack}{black!50}
\colorlet{punnetttortie}{orange!60!black}

\usetikzlibrary{positioning, decorations.pathreplacing, arrows.meta}

% --- TIKZ STYLES FOR CAT DIAGRAMS ---
\tikzset{
    % This is the main placeholder style for the cat images
    catnode/.style={
        draw, 
        rectangle, 
        rounded corners, 
        minimum height=1.5cm, 
        minimum width=2.5cm, 
        align=center, 
        font=\small\bfseries
    },
    % Allele styles (the X and Y at the top/side)
    allele/.style={font=\Large\bfseries},
    female/.style={allele, color=purple!80!black},
    male/.style={allele, color=blue!80!black},
    % --- Placeholder styles for different cats ---
    % YOU CAN REPLACE THE CONTENTS OF THESE NODES WITH YOUR IMAGES
    orange_cat/.style={
        catnode, 
        fill=orange!30, 
        text=black
    },
    black_cat/.style={
        catnode, 
        fill=black!70, 
        text=white
    },
    tortie_cat/.style={
        catnode, 
        fill=orange!50!black, % A mix for tortoiseshell
        text=white
    }
}

\usetikzlibrary{
    arrows.meta, % For nicer arrow heads (e.g., -Latex)
    positioning, % For relative node placement (e.g., above=of)
    automata     % For state diagrams, loops
}
% --- END TIKZ STYLES ---

% Common date for all slides
\date{Fall 2026}
% Counter for random transition slides generated by R
\newcounter{sssectionbg}
\setcounter{sssectionbg}{1}
\newcommand{\SSsection}[1]{
  \Ssection{#1}{FIGS/random_bg_\thesssectionbg.pdf}
  \stepcounter{sssectionbg}
}

% Random transition slide generated by R for the title page
\newcommand{\SStitlepage}{
  \titlepagewithfigure{FIGS/random_title_bg.pdf}
}

% Random transition slide generated by R for the outline page
\newcommand{\SSoutlinepage}{
  \outlinepage{FIGS/random_outline_bg.pdf}
}
\newcommand{\SSsubsection}[1]{
  \Ssubsection{#1}{FIGS/random_bg_\thesssectionbg.pdf}
  \stepcounter{sssectionbg}
}

\newcommand{\SSsubsubsection}[1]{
  \Ssubsubsection{#1}{FIGS/random_bg_\thesssectionbg.pdf}
  \stepcounter{sssectionbg}
}



\usecolortheme{orchid}
%% Listings
\usepackage{listings}
\definecolor{mygreen}{rgb}{0,0.6,0}
\definecolor{mygray}{rgb}{0.5,0.5,0.5}
\definecolor{mymauve}{rgb}{0.58,0,0.82}
\definecolor{mygold}{rgb}{1,0.843,0}
\definecolor{myblue}{rgb}{0.537,0.812,0.941}

\definecolor{mygold2}{RGB}{120,105,22}
\definecolor{mygrey2}{RGB}{50,50,50}

\definecolor{lgreen}{rgb}{0.6,0.9,.6}
\definecolor{lred}{rgb}{1,0.5,.5}

\lstloadlanguages{R}
\lstset{ %
  language=R,
  backgroundcolor=\color{black!05},   % choose the background color
  basicstyle=\footnotesize\ttfamily,        % size of fonts used for the code
  breaklines=true,                 % automatic line breaking only at whitespace
  captionpos=b,                    % sets the caption-position to bottom
  commentstyle=\color{mygreen},    % comment style
  escapeinside={\%*}{*)},          % if you want to add LaTeX within your code
  keywordstyle=\color{red},       % keyword style
  stringstyle=\color{mygold},     % string literal style
  keepspaces=true,
  columns=fullflexible,
  tabsize=4,
}
% Could also do (in lstset)
% basicstyle==\fontfamily{pcr}\footnotesize
\lstdefinelanguage{Renhanced}%
  {keywords={abbreviate,abline,abs,acos,acosh,action,add1,add,%
      aggregate,alias,Alias,alist,all,anova,any,aov,aperm,append,apply,%
      approx,approxfun,apropos,Arg,args,array,arrows,as,asin,asinh,%
      atan,atan2,atanh,attach,attr,attributes,autoload,autoloader,ave,%
      axis,backsolve,barplot,basename,besselI,besselJ,besselK,besselY,%
      beta,binomial,body,box,boxplot,break,browser,bug,builtins,bxp,by,%
      c,C,call,Call,case,cat,category,cbind,ceiling,character,char,%
      charmatch,check,chol,chol2inv,choose,chull,class,close,cm,codes,%
      coef,coefficients,co,col,colnames,colors,colours,commandArgs,%
      comment,complete,complex,conflicts,Conj,contents,contour,%
      contrasts,contr,control,helmert,contrib,convolve,cooks,coords,%
      distance,coplot,cor,cos,cosh,count,fields,cov,covratio,wt,CRAN,%
      create,crossprod,cummax,cummin,cumprod,cumsum,curve,cut,cycle,D,%
      data,dataentry,date,dbeta,dbinom,dcauchy,dchisq,de,debug,%
      debugger,Defunct,default,delay,delete,deltat,demo,de,density,%
      deparse,dependencies,Deprecated,deriv,description,detach,%
      dev2bitmap,dev,cur,deviance,off,prev,,dexp,df,dfbetas,dffits,%
      dgamma,dgeom,dget,dhyper,diag,diff,digamma,dim,dimnames,dir,%
      dirname,dlnorm,dlogis,dnbinom,dnchisq,dnorm,do,dotplot,double,%
      download,dpois,dput,drop,drop1,dsignrank,dt,dummy,dump,dunif,%
      duplicated,dweibull,dwilcox,dyn,edit,eff,effects,eigen,else,%
      emacs,end,environment,env,erase,eval,equal,evalq,example,exists,%
      exit,exp,expand,expression,External,extract,extractAIC,factor,%
      fail,family,fft,file,filled,find,fitted,fivenum,fix,floor,for,%
      For,formals,format,formatC,formula,Fortran,forwardsolve,frame,%
      frequency,ftable,ftable2table,function,gamma,Gamma,gammaCody,%
      gaussian,gc,gcinfo,gctorture,get,getenv,geterrmessage,getOption,%
      getwd,gl,glm,globalenv,gnome,GNOME,graphics,gray,grep,grey,grid,%
      gsub,hasTsp,hat,heat,help,hist,home,hsv,httpclient,I,identify,if,%
      ifelse,Im,image,\%in\%,index,influence,measures,inherits,install,%
      installed,integer,interaction,interactive,Internal,intersect,%
      inverse,invisible,IQR,is,jitter,kappa,kronecker,labels,lapply,%
      layout,lbeta,lchoose,lcm,legend,length,levels,lgamma,library,%
      licence,license,lines,list,lm,load,local,locator,log,log10,log1p,%
      log2,logical,loglin,lower,lowess,ls,lsfit,lsf,ls,machine,Machine,%
      mad,mahalanobis,make,link,margin,match,Math,matlines,mat,matplot,%
      matpoints,matrix,max,mean,median,memory,menu,merge,methods,min,%
      missing,Mod,mode,model,response,mosaicplot,mtext,mvfft,na,nan,%
      names,omit,nargs,nchar,ncol,NCOL,new,next,NextMethod,nextn,%
      nlevels,nlm,noquote,NotYetImplemented,NotYetUsed,nrow,NROW,null,%
      numeric,\%o\%,objects,offset,old,on,Ops,optim,optimise,optimize,%
      options,or,order,ordered,outer,package,packages,page,pairlist,%
      pairs,palette,panel,par,parent,parse,paste,path,pbeta,pbinom,%
      pcauchy,pchisq,pentagamma,persp,pexp,pf,pgamma,pgeom,phyper,pico,%
      pictex,piechart,Platform,plnorm,plogis,plot,pmatch,pmax,pmin,%
      pnbinom,pnchisq,pnorm,points,poisson,poly,polygon,polyroot,pos,%
      postscript,power,ppoints,ppois,predict,preplot,pretty,Primitive,%
      print,prmatrix,proc,prod,profile,proj,prompt,prop,provide,%
      psignrank,ps,pt,ptukey,punif,pweibull,pwilcox,q,qbeta,qbinom,%
      qcauchy,qchisq,qexp,qf,qgamma,qgeom,qhyper,qlnorm,qlogis,qnbinom,%
      qnchisq,qnorm,qpois,qqline,qqnorm,qqplot,qr,Q,qty,qy,qsignrank,%
      qt,qtukey,quantile,quasi,quit,qunif,quote,qweibull,qwilcox,%
      rainbow,range,rank,rbeta,rbind,rbinom,rcauchy,rchisq,Re,read,csv,%
      csv2,fwf,readline,socket,real,Recall,rect,reformulate,regexpr,%
      relevel,remove,rep,repeat,replace,replications,report,require,%
      resid,residuals,restart,return,rev,rexp,rf,rgamma,rgb,rgeom,R,%
      rhyper,rle,rlnorm,rlogis,rm,rnbinom,RNGkind,rnorm,round,row,%
      rownames,rowsum,rpois,rsignrank,rstandard,rstudent,rt,rug,runif,%
      rweibull,rwilcox,sample,sapply,save,scale,scan,scan,screen,sd,se,%
      search,searchpaths,segments,seq,sequence,setdiff,setequal,set,%
      setwd,show,sign,signif,sin,single,sinh,sink,solve,sort,source,%
      spline,splinefun,split,sqrt,stars,start,stat,stem,step,stop,%
      storage,strstrheight,stripplot,strsplit,structure,strwidth,sub,%
      subset,substitute,substr,substring,sum,summary,sunflowerplot,svd,%
      sweep,switch,symbol,symbols,symnum,sys,status,system,t,table,%
      tabulate,tan,tanh,tapply,tempfile,terms,terrain,tetragamma,text,%
      time,title,topo,trace,traceback,transform,tri,trigamma,trunc,try,%
      ts,tsp,typeof,unclass,undebug,undoc,union,unique,uniroot,unix,%
      unlink,unlist,unname,untrace,update,upper,url,UseMethod,var,%
      variable,vector,Version,vi,warning,warnings,weighted,weights,%
      which,while,window,write,\%x\%,x11,X11,xedit,xemacs,xinch,xor,%
      xpdrows,xy,xyinch,yinch,zapsmall,zip},%
   otherkeywords={!,!=,~,$,*,\%,\&,\%/\%,\%*\%,\%\%,<-,<<-,_,/},%
   alsoother={._$},%
   sensitive,%
   morecomment=[l]\#,%
   morestring=[d]",%
   morestring=[d]'% 2001 Robert Denham
  }%

% Custom color for listings
\definecolor{codegreen}{rgb}{0,0.6,0}
\definecolor{codegray}{rgb}{0.5,0.5,0.5}
\definecolor{codepurple}{rgb}{0.58,0,0.82}
\definecolor{backcolour}{rgb}{0.95,0.95,0.92}

%------------------------------------------------------------
% Listings configuration for Julia
%------------------------------------------------------------
% Define the Julia language for the listings package
\lstdefinelanguage{Julia}
{
  morekeywords={abstract,baremodule,begin,break,catch,const,continue,do,else,elseif,end,export,false,finally,for,function,global,if,import,let,local,macro,module,quote,return,struct,true,try,using,while,in},
  sensitive=true,
  comment=[l]{\#},
  morecomment=[n]{\#=}{=\#}, % For block comments like #= ... =#
  string=[b]",
  morestring=[b]',
  alsoletter={!}, % Allow ! in identifiers for mutating functions
}

% Define the style for the code blocks
\lstdefinestyle{JuliaStyle}{
    backgroundcolor=\color{backcolour},   
    commentstyle=\color{codegreen},
    keywordstyle=\color{magenta},
    numberstyle=\tiny\color{codegray},
    stringstyle=\color{codepurple},
    basicstyle=\ttfamily\footnotesize,
    breakatwhitespace=false,         
    breaklines=true,                 
    captionpos=b,                    
    keepspaces=true,                 
    % numbers=left,                    
    % numbersep=5pt,                  
    showspaces=false,                
    showstringspaces=false,
    showtabs=false,                  
    tabsize=2,
    language=Julia % Use the newly defined Julia language
}
\lstset{style=JuliaStyle}

%%%%%%% 
%% Definitions in yellow boxes
\usepackage{etoolbox}
\setbeamercolor{block title}{use=structure,fg=structure.fg,bg=structure.fg!40!bg}
\setbeamercolor{block body}{parent=normal text,use=block title,bg=block title.bg!20!bg}

% Important definition
\BeforeBeginEnvironment{importantdefinition}{%
	\setbeamercolor{block title}{fg=black,bg=yellow!20!white}
	\setbeamercolor{block body}{fg=black, bg=yellow!05!white}
}
\AfterEndEnvironment{importantdefinition}{
	\setbeamercolor{block title}{use=structure,fg=structure.fg,bg=structure.fg!20!bg}
	\setbeamercolor{block body}{parent=normal text,use=block title,bg=block title.bg!50!bg, fg=black}
}
% Regular (less important) definition
\BeforeBeginEnvironment{definition}{%
	\setbeamercolor{block title}{fg=black,bg=black!15!white}
	\setbeamercolor{block body}{fg=black, bg=black!05!white}
}
\AfterEndEnvironment{definition}{
	\setbeamercolor{block title}{use=structure,fg=structure.fg,bg=structure.fg!20!bg}
	\setbeamercolor{block body}{parent=normal text,use=block title,bg=block title.bg!50!bg, fg=black}
}
% Important theorem (red)
\BeforeBeginEnvironment{importanttheorem}{%
	\setbeamercolor{block title}{fg=black,bg=red!20!white}
	\setbeamercolor{block body}{fg=black, bg=red!05!white}
}
\AfterEndEnvironment{importanttheorem}{
	\setbeamercolor{block title}{use=structure,fg=structure.fg,bg=structure.fg!20!bg}
	\setbeamercolor{block body}{parent=normal text,use=block title,bg=block title.bg!50!bg, fg=black}
}
\BeforeBeginEnvironment{importantproperty}{%
	\setbeamercolor{block title}{fg=black,bg=red!50!white}
	\setbeamercolor{block body}{fg=black, bg=red!30!white}
}
\AfterEndEnvironment{importantproperty}{
	\setbeamercolor{block title}{use=structure,fg=structure.fg,bg=structure.fg!20!bg}
	\setbeamercolor{block body}{parent=normal text,use=block title,bg=block title.bg!50!bg, fg=black}
}

%% Colours for title page, outline and the like
\definecolor{UMbrown}{RGB}{79,44,29}
\definecolor{HighlightGold}{RGB}{230,165,83}
\definecolor{DeepGreen}{RGB}{27,46,28}
\definecolor{MediumGreen}{RGB}{52,128,56}
\definecolor{LightGreen}{RGB}{150,224,154}
\definecolor{AccentRed}{RGB}{196,44,27}
%% These are the titles used in practice
% Title page
\colorlet{title_page_title_colour}{UMbrown}
\colorlet{title_page_author_colour}{UMbrown}
\colorlet{title_page_date_colour}{UMbrown}
\colorlet{title_page_institute_colour}{UMbrown}
% Outline page
\colorlet{outline_page_title_colour}{DeepGreen}
\colorlet{outline_page_list_colour}{UMbrown}
% Section page
\colorlet{section_page_title_colour}{DeepGreen}
\colorlet{section_page_list_colour}{UMbrown}
% Subection page
\colorlet{subsection_page_title_colour}{DeepGreen}
\colorlet{subsection_page_list_colour}{UMbrown}





\IfFileExists{upquote.sty}{\usepackage{upquote}}{}
\begin{document}

% Set up cross-references and counter persistence
\setcounter{theorem}{106}

%%%%%%%%%%%%%%%%%%%%%%%%%%%%%%%%%
%%%%%%%%%%%%%%%%%%%%%%%%%%%%%%%%%
%% TITLE AND OUTLINE
%%%%%%%%%%%%%%%%%%%%%%%%%%%%%%%%%
%%%%%%%%%%%%%%%%%%%%%%%%%%%%%%%%%
\SStitlepage
\SSoutlinepage


%%%%%%%%%%%%%%%%%%%%
%%%%%%%%%%%%%%%%%%%%
%%%%%%%%%%%%%%%%%%%%
%%%%%%%%%%%%%%%%%%%%
\SSsection{Non-negative matrices and their digraphs}

\begin{frame}{Why study non-negative matrices and digraphs together?}
Many matrices in data science have entries that represent quantities that cannot be negative:
\begin{itemize}
\item counts, weights, distances or similarity scores
\item transition probabilities between states in a stochastic process
\item hyperlinks between web pages or flows in a network
\end{itemize}
\vfill
Such matrices have a rich geometric and spectral theory (Perron-Frobenius theory)
\vfill
Their zero/non-zero sparsity pattern is naturally encoded by a \defword{directed graph} (\defword{digraph}), turning algebraic questions about matrix powers into visual questions about walks on a graph
\end{frame}

\begin{frame}{Non-negative and positive vectors and matrices}
\begin{importantdefinition}[Non-negative and positive matrices]
Let $A = [a_{ij}] \in \M_{np}(\IR)$ be a real matrix
\begin{itemize}
\item $A$ is \defword{non-negative}, written $A \ge 0$ (or $A \in \M_{np}(\IRplus)$), if
\[
a_{ij} \ge 0 \quad \text{for all } i = 1,\ldots,n \text{ and } j = 1,\ldots,p
\]
\item $A$ is \defword{positive} (or \defword{strictly positive}), written $A > 0$, if
\[
a_{ij} > 0 \quad \text{for all } i = 1,\ldots,n \text{ and } j = 1,\ldots,p
\]
\end{itemize}
\end{importantdefinition}
\vfill
The same notation applies to vectors $\bx = (x_1,\ldots,x_n)^T \in \IR^n$:
\begin{itemize}
\item $\bx \ge \b0$ means $x_i \ge 0$ for all $i = 1,\ldots,n$
\item $\bx > \b0$ means $x_i > 0$ for all $i = 1,\ldots,n$
\end{itemize}
\end{frame}

\begin{frame}{Componentwise inequalities and basic properties}
For two matrices $A, B \in \M_{np}(\IR)$, we write $A \ge B$ if $A - B \ge 0$, and $A > B$ if $A - B > 0$
\vfill
\begin{proposition}[Arithmetic of non-negative matrices]
Let $A, B \in \M_n(\IRplus)$ and $\bx \in \IRplus^n$
\begin{enumerate}
\item $A + B \ge 0$ and $AB \ge 0$
\item For any integer $k \ge 1$, $A^k \ge 0$
\item If $A > 0$ and $\bx \ge \b0$ with $\bx \neq \b0$, then $A\bx > \b0$
\item If $A \ge 0$ and $\bx > \b0$, then $A\bx > \b0$ if and only if $A$ has no zero row
\end{enumerate}
\end{proposition}
\vfill
Why does property 3 hold? Because $\bx \neq \b0$ has at least one entry $x_{j_0} > 0$, so for every row $i$,
\[
(A\bx)_i = \sum_{j=1}^n a_{ij} x_j \ge a_{i,j_0} x_{j_0} > 0
\]
\end{frame}

\begin{frame}{Directed graphs (digraphs)}
\begin{importantdefinition}[Directed graph]
A \defword{directed graph} (or \defword{digraph}) is a pair $G = (V, A)$ where:
\begin{itemize}
\item $V = \{1, 2, \ldots, n\}$ is a finite set of elements called \defword{vertices} (or \defword{nodes})
\item $A \subseteq V \times V$ is a set of ordered pairs of vertices called \defword{arcs} (or \defword{directed edges})
\end{itemize}
\end{importantdefinition}
\vfill
\begin{definition}[Arc, tail, head and loop]
An \defword{arc} $e = (u, v) \in A$ represents a directed connection from vertex $u$ to vertex $v$
\begin{itemize}
\item $u$ is the \defword{initial endpoint} (or \defword{tail}) and $v$ is the \defword{terminal endpoint} (or \defword{head})
\item $u$ is a \defword{predecessor} of $v$ and $v$ is a \defword{successor} of $u$
\item An arc of the form $(u, u) \in A$ from a vertex to itself is called a \defword{loop} (or \defword{self-loop})
\end{itemize}
\end{definition}
\end{frame}

\begin{frame}{Indegree, outdegree, sources and sinks}
Let $G = (V, A)$ be a digraph and $v \in V$
\begin{definition}[Outdegree and indegree]
\begin{itemize}
\item The \defword{outdegree} $d_G^+(v)$ is the number of arcs directed away from $v$:
\[
d_G^+(v) = |\{w \in V : (v, w) \in A\}|
\]
\item The \defword{indegree} $d_G^-(v)$ is the number of arcs directed towards $v$:
\[
d_G^-(v) = |\{u \in V : (u, v) \in A\}|
\]
\item A vertex with $d_G^-(v) = 0$ is a \defword{source}; a vertex with $d_G^+(v) = 0$ is a \defword{sink}
\end{itemize}
\end{definition}
\vfill
Since every arc has one tail and one head,
\[
\sum_{v \in V} d_G^+(v) = \sum_{v \in V} d_G^-(v) = |A|
\]
\end{frame}

\begin{frame}{The digraph associated with a non-negative matrix}
\begin{importantdefinition}[Digraph of a matrix and adjacency matrix]
\begin{itemize}
\item Given a square non-negative matrix $M = [m_{ij}] \in \M_n(\IRplus)$, its associated \defword{digraph} is $\G(M) = (V, A)$ with vertex set $V = \{1, \ldots, n\}$ and arc set
\[
A = \{(i, j) \in V \times V : m_{ij} > 0\}
\]
If we attach weight $m_{ij}$ to each arc $(i, j) \in A$, we obtain a \defword{weighted digraph} (or \defword{directed network})
\item Conversely, given a digraph $G = (V, A)$ on $V = \{1, \ldots, n\}$, its (binary) \defword{adjacency matrix} is $A(G) = [a_{ij}] \in \M_n(\IRplus)$ defined by
\[
a_{ij} = \begin{cases}
1 & \text{if } (i, j) \in A \\
0 & \text{otherwise}
\end{cases}
\]
\end{itemize}
\end{importantdefinition}
\vfill
Note that diagonal entries $m_{ii} > 0$ correspond to self-loops $(i, i)$ at vertex $i$
\end{frame}

\begin{frame}{Orientation convention: general matrices vs Markov chains}
There are two natural conventions for drawing arcs from a matrix $M = [m_{ij}]$:
\vfill
\begin{enumerate}
\item \textbf{Row-to-column (adjacency) convention}: arc $(i, j)$ goes from $i$ to $j$ when $m_{ij} > 0$
\begin{itemize}
\item Row $i$ lists outgoing arcs from $i$; column $j$ lists incoming arcs into $j$
\end{itemize}
\vfill
\item \textbf{Column-to-row (transition) convention}: in column-stochastic Markov chains, $p_{ij} = \IP(S_i(t+1) \mid S_j(t))$ is the probability of moving \textbf{from $j$ to $i$}, so we draw an arc from $j$ to $i$ when $p_{ij} > 0$
\begin{itemize}
\item Column $j$ lists outgoing transitions from state $j$ (summing to $1$); row $i$ lists incoming transitions into state $i$
\item This transition digraph is simply $\G(P^T)$, obtained by reversing all arrows of $\G(P)$
\end{itemize}
\end{enumerate}
\vfill
\textbf{Crucial fact}: Reversing all arrows preserves closed walks, cycle lengths and strong connectedness, so irreducibility and primitivity of $M$ and $M^T$ are identical
\end{frame}

\begin{frame}{Example: $3 \times 3$ non-negative matrix and its digraph $\G(M)$}
Consider the $3 \times 3$ non-negative matrix $M$ and its binary sparsity pattern $A(\G(M))$:
\[
M = \begin{pmatrix}
2 & 0 & 4 \\
1 & 3 & 0 \\
0 & 5 & 0
\end{pmatrix},
\qquad
A(\G(M)) = \begin{pmatrix}
1 & 0 & 1 \\
1 & 1 & 0 \\
0 & 1 & 0
\end{pmatrix}
\]
\vfill
\begin{columns}[c]
\begin{column}{0.48\textwidth}
\centering
\textbf{Weighted digraph $\G(M)$}\\[0.3cm]
\begin{tikzpicture}[
    ->,
    >=Latex,
    every path/.style={thick},
    state_node/.style={circle, draw=UMbrown, fill=HighlightGold!25, thick, minimum size=9mm, font=\bfseries}
]
    \node[state_node] (1) at (0, 2) {$1$};
    \node[state_node] (2) at (-1.8, 0) {$2$};
    \node[state_node] (3) at (1.8, 0) {$3$};
    \path (1) edge [loop above] node[above] {$2$} (1);
    \path (1) edge [bend left=15] node[right] {$4$} (3);
    \path (2) edge [bend left=15] node[left] {$1$} (1);
    \path (2) edge [loop left] node[left] {$3$} (2);
    \path (3) edge [bend left=15] node[below] {$5$} (2);
\end{tikzpicture}
\end{column}
\begin{column}{0.52\textwidth}
Reading $\G(M)$ directly from $M$:
\begin{itemize}
\item $m_{11} = 2 > 0 \implies$ self-loop $(1, 1)$
\item $m_{13} = 4 > 0 \implies$ arc $(1, 3)$
\item $m_{21} = 1 > 0 \implies$ arc $(2, 1)$
\item $m_{22} = 3 > 0 \implies$ self-loop $(2, 2)$
\item $m_{32} = 5 > 0 \implies$ arc $(3, 2)$
\item All other entries are $0$ (no other arcs)
\end{itemize}
\end{column}
\end{columns}
\end{frame}

\begin{frame}{Example: $4 \times 4$ sparse matrix and its digraph $\G(M)$}
\begin{columns}[c]
\begin{column}{0.45\textwidth}
\[
M = \begin{pmatrix}
0 & 1 & 0 & 0 \\
0 & 0 & 1 & 1 \\
1 & 0 & 0 & 0 \\
0 & 0 & 1 & 1
\end{pmatrix}
\]
\vspace{0.2cm}
\begin{itemize}
\item Row sums of $A(\G(M))$ give the outdegrees:
\[
d^+ = (1, 2, 1, 2)^T
\]
\item Column sums of $A(\G(M))$ give the indegrees:
\[
d^- = (1, 1, 2, 2)^T
\]
\end{itemize}
\end{column}
\begin{column}{0.55\textwidth}
\centering
\begin{tikzpicture}[
    ->,
    >=Latex,
    every path/.style={thick},
    state_node/.style={circle, draw=DeepGreen, fill=LightGreen!35, thick, minimum size=9mm, font=\bfseries}
]
    \node[state_node] (1) at (0, 2.2) {$1$};
    \node[state_node] (2) at (2.8, 2.2) {$2$};
    \node[state_node] (3) at (0, 0) {$3$};
    \node[state_node] (4) at (2.8, 0) {$4$};
    \path (1) edge node[above] {$1$} (2);
    \path (2) edge node[above, sloped] {$1$} (3);
    \path (2) edge node[right] {$1$} (4);
    \path (3) edge node[left] {$1$} (1);
    \path (4) edge node[below] {$1$} (3);
    \path (4) edge [loop right] node[right] {$1$} (4);
\end{tikzpicture}
\end{column}
\end{columns}
\end{frame}

\begin{frame}[fragile]{Building a digraph from a matrix in \code{igraph}}
\begin{knitrout}
\definecolor{shadecolor}{rgb}{0.969, 0.969, 0.969}\color{fgcolor}\begin{kframe}
\begin{alltt}
\hldef{M} \hlkwb{<-} \hlkwd{matrix}\hldef{(}\hlkwd{c}\hldef{(}\hlnum{0}\hldef{,} \hlnum{1}\hldef{,} \hlnum{0}\hldef{,} \hlnum{0}\hldef{,}
              \hlnum{0}\hldef{,} \hlnum{0}\hldef{,} \hlnum{1}\hldef{,} \hlnum{1}\hldef{,}
              \hlnum{1}\hldef{,} \hlnum{0}\hldef{,} \hlnum{0}\hldef{,} \hlnum{0}\hldef{,}
              \hlnum{0}\hldef{,} \hlnum{0}\hldef{,} \hlnum{1}\hldef{,} \hlnum{1}\hldef{),} \hlkwc{nrow} \hldef{=} \hlnum{4}\hldef{,} \hlkwc{byrow} \hldef{=} \hlnum{TRUE}\hldef{)}
\hldef{G} \hlkwb{<-} \hlkwd{graph_from_adjacency_matrix}\hldef{(M,} \hlkwc{mode} \hldef{=} \hlsng{"directed"}\hldef{)}
\hlkwd{degree}\hldef{(G,} \hlkwc{mode} \hldef{=} \hlsng{"out"}\hldef{)}
\end{alltt}
\begin{verbatim}
## [1] 1 2 1 2
\end{verbatim}
\begin{alltt}
\hlkwd{degree}\hldef{(G,} \hlkwc{mode} \hldef{=} \hlsng{"in"}\hldef{)}
\end{alltt}
\begin{verbatim}
## [1] 1 1 2 2
\end{verbatim}
\begin{alltt}
\hlkwd{plot}\hldef{(G,} \hlkwc{edge.arrow.size} \hldef{=} \hlnum{0.7}\hldef{,} \hlkwc{vertex.color} \hldef{=} \hlsng{"lightblue"}\hldef{,}
     \hlkwc{vertex.size} \hldef{=} \hlnum{30}\hldef{,} \hlkwc{main} \hldef{=} \hlsng{"Digraph G(M) from adjacency matrix M"}\hldef{)}
\end{alltt}
\end{kframe}
\end{knitrout}
\end{frame}

\maxFrameImage{FIGS/U08-u08-igraph-digraph-1.pdf}


%%%%%%%%%%%%%%%%%%%%
%%%%%%%%%%%%%%%%%%%%
%%%%%%%%%%%%%%%%%%%%
%%%%%%%%%%%%%%%%%%%%
\SSsection{Matrix powers, walks and connectivity}

%%%%%%%%%%%%%%%%%%%%
%%%%%%%%%%%%%%%%%%%%
\SSsubsection{Walks and powers of the adjacency matrix}

\begin{frame}{Walks, trails, paths and cycles in a digraph}
Let $G = (V, A)$ be a digraph
\begin{importantdefinition}[Directed walk, length and closed walk]
\begin{itemize}
\item A \defword{directed walk} from vertex $v_0$ to vertex $v_k$ in $G$ is a sequence of vertices $v_0, v_1, \ldots, v_k$ such that $(v_{r-1}, v_r) \in A$ is an arc in $G$ for all $r = 1, \ldots, k$
\item The \defword{length} of the directed walk is $k$, the number of arcs traversed
\item The walk is \defword{closed} if its initial and terminal vertices coincide ($v_0 = v_k$)
\end{itemize}
\end{importantdefinition}
\vfill
\begin{definition}[Directed trail, path and cycle]
\begin{itemize}
\item A \defword{directed trail} is a directed walk in which all arcs are distinct
\item A \defword{directed path} is a directed walk in which all vertices $v_0, v_1, \ldots, v_k$ are distinct
\item A \defword{directed cycle} (of length $k \ge 1$) is a closed walk $v_0, v_1, \ldots, v_k = v_0$ in which $v_0, v_1, \ldots, v_{k-1}$ are distinct (a cycle of length $1$ is a self-loop)
\end{itemize}
\end{definition}
\end{frame}

\begin{frame}{Powers of a matrix and walks in its digraph}
\begin{importanttheorem}[Matrix powers count walks]
\label{th:matrix-powers-walks}
Let $A = [a_{ij}] \in \M_n(\IRplus)$ and denote the $(i, j)$ entry of $A^k$ by $a_{ij}^{(k)}$ for $k \in \IN_+$. Then
\[
a_{ij}^{(k)} = \sum_{r_1=1}^n \sum_{r_2=1}^n \cdots \sum_{r_{k-1}=1}^n a_{i, r_1} a_{r_1, r_2} \cdots a_{r_{k-1}, j}
\]
In particular:
\begin{enumerate}
\item $a_{ij}^{(k)} > 0$ if and only if there is at least one directed walk of length $k$ from $i$ to $j$ in $\G(A)$
\item If $A \in \{0, 1\}^{n \times n}$ is a binary adjacency matrix, then $a_{ij}^{(k)}$ equals the exact number of directed walks of length $k$ from vertex $i$ to vertex $j$ in $\G(A)$
\end{enumerate}
\end{importanttheorem}
\end{frame}

\begin{frame}[red]{Proof of Theorem~\ref{th:matrix-powers-walks}}
We prove the formula by mathematical induction on the walk length $k \ge 1$
\vfill
\textbf{Base case ($k = 1$)}: $A^1 = A$, and $a_{ij}^{(1)} = a_{ij}$ is the weight of the single arc $(i, j)$ (a walk of length $1$ from $i$ to $j$), which equals $1$ if $(i, j)$ is an arc in the binary case and $0$ otherwise
\vfill
\textbf{Inductive step}: Suppose the result holds for $k \ge 1$. Since $A^{k+1} = A^k A$, the definition of matrix multiplication gives
\[
a_{ij}^{(k+1)} = \sum_{r=1}^n a_{ir}^{(k)} a_{rj}
\]
In the binary case $A \in \{0, 1\}^{n \times n}$, any walk of length $k+1$ from $i$ to $j$ decomposes uniquely into a walk of length $k$ from $i$ to some penultimate vertex $r$, followed by the arc $(r, j)$. For a fixed $r$, there are $a_{ir}^{(k)} a_{rj}$ such walks. Summing over all $r \in \{1, \ldots, n\}$ gives $a_{ij}^{(k+1)}$, completing the induction \qed
\end{frame}

\begin{frame}{Example: counting walks of length $2$ and $3$}
Consider the binary adjacency matrix $A$ and its digraph $\G(A)$:
\begin{columns}[c]
\begin{column}{0.45\textwidth}
\centering
\begin{tikzpicture}[
    ->,
    >=Latex,
    every path/.style={thick},
    state_node/.style={circle, draw=UMbrown, fill=HighlightGold!25, thick, minimum size=9mm, font=\bfseries}
]
    \node[state_node] (1) at (0, 1.8) {$1$};
    \node[state_node] (2) at (-1.6, 0) {$2$};
    \node[state_node] (3) at (1.6, 0) {$3$};
    \path (1) edge [loop above] (1);
    \path (1) edge [bend left=15] (2);
    \path (2) edge [bend left=15] (1);
    \path (2) edge (3);
    \path (3) edge (1);
\end{tikzpicture}
\end{column}
\begin{column}{0.55\textwidth}
\[
A = \begin{pmatrix}
1 & 1 & 0 \\
1 & 0 & 1 \\
1 & 0 & 0
\end{pmatrix},
\quad
A^2 = \begin{pmatrix}
2 & 1 & 1 \\
2 & 1 & 0 \\
1 & 1 & 0
\end{pmatrix}
\]
\[
A^3 = A^2 A = \begin{pmatrix}
4 & 2 & 1 \\
3 & 2 & 1 \\
2 & 1 & 1
\end{pmatrix}
\]
\end{column}
\end{columns}
\vfill
Let us verify individual entries directly from $\G(A)$:
\begin{itemize}
\item $(A^2)_{11} = 2$: the two walks of length $2$ from $1$ to $1$ are $1 \to 1 \to 1$ and $1 \to 2 \to 1$
\item $(A^2)_{23} = 0$: from $2$, length-$1$ steps reach $1$ or $3$, neither of which has an arc to $3$
\item $(A^3)_{11} = 4$: the four walks of length $3$ from $1$ to $1$ are
\[
1 \to 1 \to 1 \to 1, \quad 1 \to 1 \to 2 \to 1, \quad 1 \to 2 \to 1 \to 1, \quad 1 \to 2 \to 3 \to 1
\]
\end{itemize}
\end{frame}


%%%%%%%%%%%%%%%%%%%%
%%%%%%%%%%%%%%%%%%%%
\SSsubsection{Irreducibility and strong connectedness}

\begin{frame}{Strong connectedness and strongly connected components}
\begin{importantdefinition}[Strongly connected digraph]
A digraph $G = (V, A)$ is \defword{strongly connected} if for every ordered pair of vertices $(i, j) \in V \times V$ with $i \neq j$, there exists a directed walk (and hence a directed path) in $G$ from $i$ to $j$ and from $j$ to $i$
\end{importantdefinition}
\vfill
Define the mutual reachability relation $i \equiv j$ on $V$ by:
\[
i \equiv j \iff i = j \text{ or there is a directed walk from } i \text{ to } j \text{ and from } j \text{ to } i
\]
\begin{itemize}
\item $\equiv$ is reflexive, symmetric and transitive, so it is an \defword{equivalence relation} on $V$
\item Its equivalence classes partition $V$ into maximal strongly connected subdigraphs called the \defword{strongly connected components} (SCCs) of $G$
\item $G$ is strongly connected if and only if it has a single strongly connected component
\end{itemize}
\end{frame}

\begin{frame}{Reducible and irreducible matrices}
\begin{importantdefinition}[Reducible and irreducible matrices]
A matrix $A \in \M_n(\IRplus)$ ($n \ge 2$) is \defword{reducible} if there exists a permutation matrix $P_\pi \in \M_n$ such that
\[
P_\pi^T A P_\pi = \begin{pmatrix}
A_{11} & A_{12} \\
\b0 & A_{22}
\end{pmatrix}
\]
where $A_{11} \in \M_r$ and $A_{22} \in \M_{n-r}$ are square blocks with $1 \le r < n$, and $\b0$ is the $(n-r) \times r$ zero matrix
\vskip0.2cm
The matrix $A$ is \defword{irreducible} if no such permutation matrix exists (by convention, any $1 \times 1$ matrix is irreducible)
\end{importantdefinition}
\vfill
Note that $P_\pi^T A P_\pi$ simply relabels the rows and columns of $A$ using the same permutation $\pi$ of the vertices $\{1, \ldots, n\}$
\end{frame}

\begin{frame}{Why block upper-triangular form prevents strong connectedness}
Suppose $A$ has been relabelled into the reducible block form
\[
\tilde{A} = P_\pi^T A P_\pi = \begin{pmatrix}
A_{11} & A_{12} \\
\b0 & A_{22}
\end{pmatrix}
\begin{array}{l}
\}\; V_1 = \{1, \ldots, r\} \\
\}\; V_2 = \{r+1, \ldots, n\}
\end{array}
\]
\vfill
\begin{itemize}
\item The bottom-left block is $\b0$, which means $\tilde{a}_{ij} = 0$ for all $i \in V_2$ and $j \in V_1$: there is \textbf{no arc} from any vertex in $V_2$ to any vertex in $V_1$
\item Furthermore, for every power $k \ge 1$, block multiplication yields
\[
\tilde{A}^k = \begin{pmatrix}
A_{11}^k & * \\
\b0 & A_{22}^k
\end{pmatrix}
\]
so $(\tilde{A}^k)_{ij} = 0$ for all $i \in V_2$ and $j \in V_1$
\item Thus no walk of any length $k \ge 1$ can ever go from $V_2$ to $V_1$: the vertex set $V_2$ is a \defword{closed trap} under the row-to-column walk dynamics
\end{itemize}
\end{frame}

\begin{frame}{Characterisations of irreducibility}
\begin{importanttheorem}[Characterisations of irreducible matrices]
\label{th:irreducible-tfae}
For a non-negative matrix $A \in \M_n(\IRplus)$, the following statements are equivalent:
\begin{enumerate}
\item $A$ is irreducible
\item The associated digraph $\G(A)$ is strongly connected
\item $\II_n + A + A^2 + \cdots + A^{n-1} > 0$ (strictly positive entrywise)
\item $(\II_n + A)^{n-1} > 0$ (strictly positive entrywise)
\end{enumerate}
\end{importanttheorem}
\vfill
Items 3 and 4 provide a purely algebraic test for irreducibility without searching over all $n!$ permutation matrices
\end{frame}

\begin{frame}[red]{Proof of Theorem~\ref{th:irreducible-tfae}}
\textbf{(1) $\iff$ (2)}: $A$ is reducible $\iff$ vertices can be partitioned into non-empty sets $V_1, V_2$ with no arc from $V_2$ to $V_1$ $\iff$ no walk can go from $V_2$ to $V_1$ $\iff$ $\G(A)$ is not strongly connected
\vfill
\textbf{(2) $\iff$ (3)}: By Theorem~\ref{th:matrix-powers-walks}, $(A^k)_{ij} > 0 \iff$ there is a walk of length $k$ from $i$ to $j$. Any shortest walk between two distinct vertices $i \neq j$ is a simple path visiting at most $n$ vertices, so its length is $k \in \{1, \ldots, n-1\}$. With $(\II_n)_{ii} = 1 > 0$ covering $i = j$, $\G(A)$ is strongly connected $\iff$ for every pair $(i, j)$ there is some $k \in \{0, 1, \ldots, n-1\}$ with $(A^k)_{ij} > 0$ $\iff$ $\sum_{k=0}^{n-1} A^k > 0$
\vfill
\textbf{(3) $\iff$ (4)}: By the binomial expansion with positive coefficients $\binom{n-1}{k} > 0$ and $A^k \ge 0$,
\[
(\II_n + A)^{n-1} = \sum_{k=0}^{n-1} \binom{n-1}{k} A^k
\]
A sum of non-negative terms with positive weights is strictly positive if and only if at least one term is strictly positive, so $(\II_n + A)^{n-1} > 0 \iff \sum_{k=0}^{n-1} A^k > 0$ \qed
\end{frame}

\begin{frame}{Exam-style example 1: testing irreducibility via $\G(A)$}
\textbf{Problem}: Determine whether the following $4 \times 4$ matrix is irreducible:
\[
A = \begin{pmatrix}
0 & 2 & 0 & 0 \\
0 & 0 & 1 & 3 \\
0 & 0 & 0 & 1 \\
4 & 0 & 0 & 0
\end{pmatrix}
\]
\begin{columns}[c]
\begin{column}{0.45\textwidth}
\centering
\begin{tikzpicture}[
    ->,
    >=Latex,
    every path/.style={thick},
    state_node/.style={circle, draw=UMbrown, fill=HighlightGold!25, thick, minimum size=8.5mm, font=\bfseries}
]
    \node[state_node] (1) at (0, 1.8) {$1$};
    \node[state_node] (2) at (2.4, 1.8) {$2$};
    \node[state_node] (3) at (2.4, 0) {$3$};
    \node[state_node] (4) at (0, 0) {$4$};
    \path (1) edge node[above] {$2$} (2);
    \path (2) edge node[right] {$1$} (3);
    \path (2) edge node[above, sloped] {$3$} (4);
    \path (3) edge node[below] {$1$} (4);
    \path (4) edge node[left] {$4$} (1);
\end{tikzpicture}
\end{column}
\begin{column}{0.55\textwidth}
\textbf{Solution}:
\begin{itemize}
\item Draw the 4-node digraph $\G(A)$
\item Observe the directed Hamiltonian cycle
\[
1 \to 2 \to 3 \to 4 \to 1
\]
\item Since this cycle visits all 4 vertices, every vertex can reach every other vertex along the cycle
\item Therefore $\G(A)$ is \textbf{strongly connected}, so $A$ is \textbf{irreducible}
\end{itemize}
\end{column}
\end{columns}
\end{frame}

\begin{frame}{Exam-style example 2: reducible matrix and block triangular form}
\textbf{Problem}: Show that $B$ is reducible, find its SCCs and write $P_\pi^T B P_\pi$ in block upper-triangular form:
\[
B = \begin{pmatrix}
0 & 0 & 1 & 0 \\
1 & 0 & 0 & 1 \\
1 & 0 & 0 & 0 \\
0 & 1 & 1 & 0
\end{pmatrix}
\]
\begin{columns}[c]
\begin{column}{0.42\textwidth}
\centering
\begin{tikzpicture}[
    ->,
    >=Latex,
    every path/.style={thick},
    scc1/.style={circle, draw=DeepGreen, fill=LightGreen!40, thick, minimum size=8mm, font=\bfseries},
    scc2/.style={circle, draw=AccentRed, fill=AccentRed!20, thick, minimum size=8mm, font=\bfseries}
]
    \node[scc2] (2) at (0, 1.8) {$2$};
    \node[scc2] (4) at (0, 0) {$4$};
    \node[scc1] (1) at (2.4, 1.8) {$1$};
    \node[scc1] (3) at (2.4, 0) {$3$};
    \path (2) edge [bend left=15] (4);
    \path (4) edge [bend left=15] (2);
    \path (2) edge (1);
    \path (4) edge (3);
    \path (1) edge [bend left=15] (3);
    \path (3) edge [bend left=15] (1);
\end{tikzpicture}
\end{column}
\begin{column}{0.58\textwidth}
\begin{itemize}
\item SCCs are $C_1 = \{2, 4\}$ and $C_2 = \{1, 3\}$
\item Arcs go from $C_1$ to $C_2$, but \textbf{no arc goes from $C_2$ to $C_1$}, so $\G(B)$ is not strongly connected ($B$ is \textbf{reducible})
\item Ordering vertices as $(2, 4, 1, 3)$ gives
\[
P_\pi^T B P_\pi = \left(\begin{array}{cc|cc}
0 & 1 & 1 & 0 \\
1 & 0 & 0 & 1 \\ \hline
0 & 0 & 0 & 1 \\
0 & 0 & 1 & 0
\end{array}\right) = \begin{pmatrix} B_{11} & B_{12} \\ \b0 & B_{22} \end{pmatrix}
\]
\end{itemize}
\end{column}
\end{columns}
\end{frame}

\begin{frame}[fragile]{Checking strong connectedness and $(\II_n + A)^{n-1}$ in R}
\begin{knitrout}
\definecolor{shadecolor}{rgb}{0.969, 0.969, 0.969}\color{fgcolor}\begin{kframe}
\begin{alltt}
\hldef{A} \hlkwb{<-} \hlkwd{matrix}\hldef{(}\hlkwd{c}\hldef{(}\hlnum{0}\hldef{,} \hlnum{2}\hldef{,} \hlnum{0}\hldef{,} \hlnum{0}\hldef{,}
              \hlnum{0}\hldef{,} \hlnum{0}\hldef{,} \hlnum{1}\hldef{,} \hlnum{3}\hldef{,}
              \hlnum{0}\hldef{,} \hlnum{0}\hldef{,} \hlnum{0}\hldef{,} \hlnum{1}\hldef{,}
              \hlnum{4}\hldef{,} \hlnum{0}\hldef{,} \hlnum{0}\hldef{,} \hlnum{0}\hldef{),} \hlkwc{nrow} \hldef{=} \hlnum{4}\hldef{,} \hlkwc{byrow} \hldef{=} \hlnum{TRUE}\hldef{)}
\hlkwd{is_connected}\hldef{(}\hlkwd{graph_from_adjacency_matrix}\hldef{(A} \hlopt{>} \hlnum{0}\hldef{),} \hlkwc{mode} \hldef{=} \hlsng{"strong"}\hldef{)}
\end{alltt}
\begin{verbatim}
## [1] TRUE
\end{verbatim}
\begin{alltt}
\hlcom{# Compute (I_4 + A)^3}
\hldef{I4_plus_A} \hlkwb{<-} \hlkwd{diag}\hldef{(}\hlnum{4}\hldef{)} \hlopt{+} \hldef{A}
\hldef{I4_plus_A} \hlopt{%*%} \hldef{I4_plus_A} \hlopt{%*%} \hldef{I4_plus_A}
\end{alltt}
\begin{verbatim}
##      [,1] [,2] [,3] [,4]
## [1,]   25    6    6   20
## [2,]   40   25    3   12
## [3,]   12    8    1    3
## [4,]   12   24    8   25
\end{verbatim}
\end{kframe}
\end{knitrout}
\end{frame}


%%%%%%%%%%%%%%%%%%%%
%%%%%%%%%%%%%%%%%%%%
\SSsubsection{Primitivity and aperiodicity}

\begin{frame}{Primitive matrices}
\begin{importantdefinition}[Primitive matrix]
A non-negative matrix $A \in \M_n(\IRplus)$ is \defword{primitive} if there exists a positive integer $m \in \IN_+$ such that
\[
A^m > 0 \quad (\text{every entry of } A^m \text{ is strictly positive})
\]
The smallest positive integer $m$ such that $A^m > 0$ is called the \defword{index of primitivity} of $A$, denoted $\gamma(A)$
\vskip0.2cm
An irreducible non-negative matrix that is not primitive is called \defword{imprimitive} (or \defword{periodic})
\end{importantdefinition}
\vfill
\begin{proposition}
\begin{enumerate}
\item Every primitive matrix is irreducible
\item If $A^m > 0$ for some $m \in \IN_+$, then $A^k > 0$ for all $k \ge m$
\end{enumerate}
\end{proposition}
\end{frame}

\begin{frame}{Irreducible does not imply primitive!}
Consider the $2 \times 2$ and $3 \times 3$ cyclic permutation matrices:
\[
C_2 = \begin{pmatrix} 0 & 1 \\ 1 & 0 \end{pmatrix},
\qquad
C_3 = \begin{pmatrix} 0 & 1 & 0 \\ 0 & 0 & 1 \\ 1 & 0 & 0 \end{pmatrix}
\]
\begin{columns}[c]
\begin{column}{0.4\textwidth}
\centering
\begin{tikzpicture}[
    ->,
    >=Latex,
    every path/.style={thick},
    state_node/.style={circle, draw=UMbrown, fill=HighlightGold!25, thick, minimum size=8.5mm, font=\bfseries}
]
    \node[state_node] (1) at (0, 1.6) {$1$};
    \node[state_node] (2) at (1.8, 0.2) {$2$};
    \node[state_node] (3) at (-1.8, 0.2) {$3$};
    \path (1) edge [bend left=20] (2);
    \path (2) edge [bend left=20] (3);
    \path (3) edge [bend left=20] (1);
\end{tikzpicture}
\end{column}
\begin{column}{0.6\textwidth}
\begin{itemize}
\item Both $\G(C_2)$ and $\G(C_3)$ are single directed cycles, so $C_2$ and $C_3$ are \textbf{irreducible}
\item Yet for $C_2$: $C_2^{2k} = \II_2$ and $C_2^{2k+1} = C_2$
\item And for $C_3$: $C_3^2 = \begin{pmatrix} 0 & 0 & 1 \\ 1 & 0 & 0 \\ 0 & 1 & 0 \end{pmatrix}$, $C_3^3 = \II_3$, $C_3^4 = C_3, \ldots$
\item Every power $C_2^m$ and $C_3^m$ contains zeros, so neither matrix is primitive!
\end{itemize}
\end{column}
\end{columns}
\vfill
Why? Because walks from $i$ to $j$ only occur at lengths $m$ belonging to a fixed congruence class modulo $2$ (for $C_2$) or modulo $3$ (for $C_3$)
\end{frame}

\begin{frame}{Period of a digraph and characterisation of primitivity}
\begin{importantdefinition}[Period and aperiodicity]
Let $G = (V, A)$ be a strongly connected digraph. The \defword{period} (or \defword{index of imprimitivity}) $d(G)$ of $G$ is the greatest common divisor of the lengths of all directed cycles (or closed walks) in $G$:
\[
d(G) = \gcd\{\text{length}(C) : C \text{ is a directed cycle in } G\}
\]
\begin{itemize}
\item If $d(G) = 1$, we say that $G$ (and its matrix $A$) is \defword{aperiodic}
\item If $d(G) > 1$, we say that $G$ (and its matrix $A$) is \defword{periodic} with period $d(G)$
\end{itemize}
\end{importantdefinition}
\vfill
\begin{importanttheorem}[Graphical characterisation of primitive matrices]
\label{th:primitive-iff-sc-aperiodic}
A non-negative matrix $A \in \M_n(\IRplus)$ is \textbf{primitive} if and only if $\G(A)$ is \textbf{strongly connected} and \textbf{aperiodic} ($d(\G(A)) = 1$)
\end{importanttheorem}
\end{frame}

\begin{frame}{Practical corollaries for exams}
Checking $\gcd$ of cycle lengths is immediate in two common situations:
\vfill
\begin{corollary}[Self-loop test for primitivity]
\label{cor:self-loop-primitive}
If $A \in \M_n(\IRplus)$ is irreducible ($\G(A)$ is strongly connected) and has \textbf{at least one positive diagonal entry} $a_{ii} > 0$ for some $i \in \{1, \ldots, n\}$, then $A$ is \textbf{primitive}
\end{corollary}
\textbf{Reason}: $a_{ii} > 0$ gives a self-loop $i \to i$ of length $1$ in $\G(A)$, and $\gcd(1, \ldots) = 1$
\vfill
\begin{corollary}[Coprime cycle lengths]
If $\G(A)$ is strongly connected and contains two directed cycles of coprime lengths $\ell_1$ and $\ell_2$ (i.e. $\gcd(\ell_1, \ell_2) = 1$, such as a $2$-cycle and a $3$-cycle), then $A$ is \textbf{primitive}
\end{corollary}
\end{frame}

\begin{frame}{Visualising periodic vs primitive digraphs}
\begin{columns}[T]
\begin{column}{0.32\textwidth}
\centering
\textbf{Case 1: Irreducible, periodic ($d = 3$)}\\[0.2cm]
\begin{tikzpicture}[
    ->, >=Latex, every path/.style={thick},
    state_node/.style={circle, draw=AccentRed, fill=AccentRed!15, thick, minimum size=7.5mm, font=\bfseries}
]
    \node[state_node] (1) at (0, 1.4) {$1$};
    \node[state_node] (2) at (1.2, 0) {$2$};
    \node[state_node] (3) at (-1.2, 0) {$3$};
    \path (1) edge [bend left=18] (2);
    \path (2) edge [bend left=18] (3);
    \path (3) edge [bend left=18] (1);
\end{tikzpicture}
\\[0.2cm]
\[
A_1 = \begin{pmatrix} 0 & 1 & 0 \\ 0 & 0 & 1 \\ 1 & 0 & 0 \end{pmatrix}
\]
Cycle length: $\{3\}$\\
$\gcd(3) = 3 \implies$ \textbf{not primitive}
\end{column}
\begin{column}{0.32\textwidth}
\centering
\textbf{Case 2: Adding one self-loop ($d = 1$)}\\[0.2cm]
\begin{tikzpicture}[
    ->, >=Latex, every path/.style={thick},
    state_node/.style={circle, draw=DeepGreen, fill=LightGreen!35, thick, minimum size=7.5mm, font=\bfseries}
]
    \node[state_node] (1) at (0, 1.4) {$1$};
    \node[state_node] (2) at (1.2, 0) {$2$};
    \node[state_node] (3) at (-1.2, 0) {$3$};
    \path (1) edge [loop above] (1);
    \path (1) edge [bend left=18] (2);
    \path (2) edge [bend left=18] (3);
    \path (3) edge [bend left=18] (1);
\end{tikzpicture}
\\[0.2cm]
\[
A_2 = \begin{pmatrix} 1 & 1 & 0 \\ 0 & 0 & 1 \\ 1 & 0 & 0 \end{pmatrix}
\]
Cycle lengths: $\{1, 3\}$\\
$\gcd(1, 3) = 1 \implies$ \textbf{primitive} ($A_2^4 > 0$)
\end{column}
\begin{column}{0.32\textwidth}
\centering
\textbf{Case 3: Zero trace, coprime cycles ($d = 1$)}\\[0.2cm]
\begin{tikzpicture}[
    ->, >=Latex, every path/.style={thick},
    state_node/.style={circle, draw=DeepGreen, fill=LightGreen!35, thick, minimum size=7.5mm, font=\bfseries}
]
    \node[state_node] (1) at (0, 1.4) {$1$};
    \node[state_node] (2) at (1.2, 0) {$2$};
    \node[state_node] (3) at (-1.2, 0) {$3$};
    \path (1) edge [bend left=18] (2);
    \path (2) edge [bend left=18] (1);
    \path (2) edge [bend left=18] (3);
    \path (3) edge [bend left=18] (1);
\end{tikzpicture}
\\[0.2cm]
\[
A_3 = \begin{pmatrix} 0 & 1 & 0 \\ 1 & 0 & 1 \\ 1 & 0 & 0 \end{pmatrix}
\]
Cycles: $1 \leftrightarrow 2$ ($2$), $1 \to 2 \to 3 \to 1$ ($3$)\\
$\gcd(2, 3) = 1 \implies$ \textbf{primitive} ($A_3^5 > 0$)
\end{column}
\end{columns}
\end{frame}

\begin{frame}{Exam-style comparison: two $4 \times 4$ irreducible matrices}
Both matrices below are irreducible with zero diagonal ($\tr = 0$). Which one is primitive?
\begin{columns}[T]
\begin{column}{0.48\textwidth}
\centering
\[
M_1 = \begin{pmatrix}
0 & 1 & 0 & 0 \\
1 & 0 & 1 & 0 \\
0 & 0 & 0 & 1 \\
1 & 0 & 0 & 0
\end{pmatrix}
\]
\begin{tikzpicture}[
    ->, >=Latex, every path/.style={thick},
    state_node/.style={circle, draw=AccentRed, fill=AccentRed!15, thick, minimum size=7.5mm, font=\bfseries}
]
    \node[state_node] (1) at (0, 1.2) {$1$};
    \node[state_node] (2) at (2.0, 1.2) {$2$};
    \node[state_node] (3) at (2.0, 0) {$3$};
    \node[state_node] (4) at (0, 0) {$4$};
    \path (1) edge [bend left=15] (2);
    \path (2) edge [bend left=15] (1);
    \path (2) edge (3);
    \path (3) edge (4);
    \path (4) edge (1);
\end{tikzpicture}
\begin{itemize}
\item Directed cycles: $1 \to 2 \to 1$ (length $2$) and $1 \to 2 \to 3 \to 4 \to 1$ (length $4$)
\item $d = \gcd(2, 4) = 2 > 1$: \textbf{imprimitive} (periodic with period $2$)
\end{itemize}
\end{column}
\begin{column}{0.48\textwidth}
\centering
\[
M_2 = \begin{pmatrix}
0 & 1 & 0 & 0 \\
0 & 0 & 1 & 1 \\
0 & 0 & 0 & 1 \\
1 & 0 & 0 & 0
\end{pmatrix}
\]
\begin{tikzpicture}[
    ->, >=Latex, every path/.style={thick},
    state_node/.style={circle, draw=DeepGreen, fill=LightGreen!35, thick, minimum size=7.5mm, font=\bfseries}
]
    \node[state_node] (1) at (0, 1.2) {$1$};
    \node[state_node] (2) at (2.0, 1.2) {$2$};
    \node[state_node] (3) at (2.0, 0) {$3$};
    \node[state_node] (4) at (0, 0) {$4$};
    \path (1) edge (2);
    \path (2) edge (3);
    \path (2) edge (4);
    \path (3) edge (4);
    \path (4) edge (1);
\end{tikzpicture}
\begin{itemize}
\item Directed cycles: $1 \to 2 \to 4 \to 1$ (length $3$) and $1 \to 2 \to 3 \to 4 \to 1$ (length $4$)
\item $d = \gcd(3, 4) = 1$: \textbf{primitive} ($M_2^{10} > 0$)
\end{itemize}
\end{column}
\end{columns}
\end{frame}

\begin{frame}{Wielandt's bound on the index of primitivity}
How large can the smallest exponent $\gamma(A)$ with $A^{\gamma(A)} > 0$ be for an $n \times n$ primitive matrix?
\vfill
\begin{importanttheorem}[Wielandt's theorem]
\label{th:wielandt}
Let $A \in \M_n(\IRplus)$ with $n \ge 2$. Then $A$ is primitive if and only if
\[
A^{n^2 - 2n + 2} > 0
\]
Moreover, if $A$ is irreducible and has $d \ge 1$ positive diagonal entries, then
\[
\gamma(A) \le 2n - d - 1
\]
\end{importanttheorem}
\vfill
\begin{itemize}
\item For $n = 3$, $n^2 - 2n + 2 = 5$: the matrix $A_3$ has $A_3^4 \not> 0$ and $A_3^5 > 0$
\item For $n = 4$, $n^2 - 2n + 2 = 10$: the matrix $M_2$ on the previous slide has $M_2^9 \not> 0$ and $M_2^{10} > 0$, showing that Wielandt's bound $n^2 - 2n + 2$ is sharp!
\end{itemize}
\end{frame}


%%%%%%%%%%%%%%%%%%%%
%%%%%%%%%%%%%%%%%%%%
%%%%%%%%%%%%%%%%%%%%
%%%%%%%%%%%%%%%%%%%%
\SSsection{The Perron-Frobenius theorem}

\begin{frame}{Spectrum and spectral radius}
\begin{importantdefinition}[Spectrum and spectral radius]
Let $A \in \M_n(\IC)$ be a square matrix
\begin{itemize}
\item The \defword{spectrum} of $A$, denoted $\sigma(A)$, is the set of all eigenvalues of $A$:
\[
\sigma(A) = \{\lambda \in \IC : \det(A - \lambda \II_n) = 0\}
\]
\item The \defword{spectral radius} of $A$, denoted $\rho(A)$, is the maximum modulus of its eigenvalues:
\[
\rho(A) = \max_{\lambda \in \sigma(A)} |\lambda|
\]
\end{itemize}
\end{importantdefinition}
\vfill
Geometrically, in the complex plane $\IC$, all eigenvalues of $A$ lie inside or on the closed disc $\{z \in \IC : |z| \le \rho(A)\}$ of radius $\rho(A)$ centred at the origin
\end{frame}

\begin{frame}{Geometric intuition in $\IR^2$: trapping an invariant ray}
Why must a positive matrix $A > 0$ have a positive eigenvalue with a positive eigenvector?
\begin{columns}[c]
\begin{column}{0.48\textwidth}
\centering
\begin{tikzpicture}[scale=1.05, >=Latex]
    % Shaded positive quadrant
    \fill[HighlightGold!15] (0,0) rectangle (3.4,2.8);
    % Shaded image cone A(R_+^2)
    \fill[MediumGreen!30] (0,0) -- (3.4, 0.9) -- (3.4, 2.8) -- (1.1, 2.8) -- cycle;
    % Axes
    \draw[->, thick] (-0.3,0) -- (3.6,0) node[right] {$x_1$};
    \draw[->, thick] (0,-0.3) -- (0,3.0) node[above] {$x_2$};
    % Standard basis images A e_1 and A e_2
    \draw[->, very thick, DeepGreen] (0,0) -- (3.1, 0.82) node[below right] {$A\be_1$};
    \draw[->, very thick, DeepGreen] (0,0) -- (0.95, 2.45) node[above left] {$A\be_2$};
    % Perron ray
    \draw[->, ultra thick, AccentRed] (0,0) -- (2.8, 1.95) node[above right] {$\bx_P$ ($A\bx_P = \rho(A)\bx_P$)};
    % Simplex segment
    \draw[thick, dashed, UMbrown] (2.2,0) -- (0,2.2);
\end{tikzpicture}
\end{column}
\begin{column}{0.52\textwidth}
\begin{itemize}
\item Consider the non-negative quadrant $\IRplus^2 = \{\bx \ge \b0\}$
\item Since $A > 0$, both columns $A\be_1 > \b0$ and $A\be_2 > \b0$ lie \textbf{strictly inside} the open first quadrant
\item Thus $A$ maps the entire quadrant $\IRplus^2 \setminus \{\b0\}$ into a narrower cone strictly inside the first quadrant
\item By Brouwer's fixed-point theorem (or the intermediate value theorem on the angle of $\bx$), $A$ must map a unique ray in the first quadrant onto itself!
\end{itemize}
\end{column}
\end{columns}
\end{frame}

\begin{frame}{Perron's theorem for positive matrices}
\begin{importanttheorem}[Perron's theorem, 1907]
\label{th:perron}
Let $A \in \M_n(\IR)$ be a strictly positive matrix ($A > 0$). Then:
\begin{enumerate}
\item $\rho(A) > 0$ is an eigenvalue of $A$ (the \defword{Perron root} or \defword{dominant eigenvalue})
\item $\rho(A)$ is algebraically and geometrically simple (multiplicity $1$)
\item $\rho(A)$ strictly dominates every other eigenvalue: $|\lambda| < \rho(A)$ for all $\lambda \in \sigma(A) \setminus \{\rho(A)\}$
\item There exists a \textbf{unique} strictly positive right eigenvector $\bx > \b0$ such that
\[
A\bx = \rho(A)\bx \quad \text{and} \quad \11^T\bx = \sum_{i=1}^n x_i = 1
\]
(the \defword{right Perron vector}, and similarly a unique left Perron vector $\by > \b0$)
\item No other eigenvalue of $A$ has a non-negative eigenvector
\end{enumerate}
\end{importanttheorem}
\end{frame}

\begin{frame}{The Perron-Frobenius theorem for primitive and irreducible matrices}
Frobenius (1912) extended Perron's theorem to irreducible non-negative matrices:
\begin{importanttheorem}[Perron-Frobenius theorem]
\label{th:perron-frobenius}
Let $A \in \M_n(\IRplus)$ be an \textbf{irreducible} non-negative matrix
\begin{enumerate}
\item $\rho(A) > 0$ is a simple eigenvalue of $A$ (algebraic and geometric multiplicity $1$)
\item There exists a unique strictly positive right eigenvector $\bx > \b0$ with $A\bx = \rho(A)\bx$ and $\11^T\bx = 1$ (and no other eigenvalue has a non-negative eigenvector)
\item \textbf{If $A$ is primitive} ($d = 1$): $|\lambda| < \rho(A)$ for all $\lambda \in \sigma(A)$ with $\lambda \neq \rho(A)$
\item \textbf{If $A$ is imprimitive with period $d > 1$}: $A$ has exactly $d$ eigenvalues of modulus $\rho(A)$, given by $\lambda_k = \rho(A) e^{2\pi i k / d}$ ($k = 0, 1, \ldots, d-1$), and $|\lambda| < \rho(A)$ for all others
\end{enumerate}
\end{importanttheorem}
\end{frame}

\begin{frame}{Spectral circle: primitive ($d = 1$) vs periodic ($d = 3$)}
\begin{columns}[c]
\begin{column}{0.48\textwidth}
\centering
\textbf{Primitive matrix ($d = 1$)}\\[0.2cm]
\begin{tikzpicture}[scale=1.1, >=Latex]
    \draw[dashed, thick, UMbrown] (0,0) circle (1.6cm);
    \draw[->, thin, gray] (-2.1,0) -- (2.3,0) node[right, black] {$\Re(\lambda)$};
    \draw[->, thin, gray] (0,-2.0) -- (0,2.0) node[above, black] {$\Im(\lambda)$};
    % Dominant eigenvalue
    \fill[AccentRed] (1.6, 0) circle (3pt) node[above right] {$\rho(A)$};
    % Other eigenvalues strictly inside
    \fill[DeepGreen] (-0.7, 0.6) circle (2.5pt);
    \fill[DeepGreen] (-0.7, -0.6) circle (2.5pt);
    \fill[DeepGreen] (0.4, 0) circle (2.5pt);
\end{tikzpicture}
\\[0.2cm]
Strict spectral gap $|\lambda_j| < \rho(A)$ for $j \ge 2$\\
$\implies \left(\frac{1}{\rho(A)} A\right)^k$ \textbf{converges} as $k \to \infty$
\end{column}
\begin{column}{0.48\textwidth}
\centering
\textbf{Irreducible periodic matrix ($d = 3$)}\\[0.2cm]
\begin{tikzpicture}[scale=1.1, >=Latex]
    \draw[dashed, thick, UMbrown] (0,0) circle (1.6cm);
    \draw[->, thin, gray] (-2.1,0) -- (2.3,0) node[right, black] {$\Re(\lambda)$};
    \draw[->, thin, gray] (0,-2.0) -- (0,2.0) node[above, black] {$\Im(\lambda)$};
    % 3 eigenvalues on spectral circle
    \fill[AccentRed] (1.6, 0) circle (3pt) node[above right] {$\rho(A)$};
    \fill[AccentRed] (-0.8, 1.386) circle (3pt) node[above left] {$\rho(A)e^{2\pi i/3}$};
    \fill[AccentRed] (-0.8, -1.386) circle (3pt) node[below left] {$\rho(A)e^{4\pi i/3}$};
    % Interior eigenvalue
    \fill[DeepGreen] (0.3, 0) circle (2.5pt);
\end{tikzpicture}
\\[0.2cm]
$d = 3$ eigenvalues on the circle $|\lambda| = \rho(A)$\\
$\implies \left(\frac{1}{\rho(A)} A\right)^k$ \textbf{oscillates} with period $3$
\end{column}
\end{columns}
\end{frame}

\begin{frame}{Summary of Perron-Frobenius across all non-negative matrices}
Let $A \in \M_n(\IRplus)$ have Perron root $\lambda_P = \rho(A)$ and right/left Perron vectors $\bx_P, \by_P$:
\vfill
\begin{center}
\begin{tikzpicture}[
    scale=0.78, every node/.style={transform shape},
    box/.style={draw=UMbrown, thick, rounded corners=3pt, fill=HighlightGold!12, align=left, inner sep=5pt}
]
    \node[box, fill=UMbrown!15, font=\bfseries] (nonneg) at (0, 0) {Non-negative matrix $A \ge 0$};
    \node[box, text width=4.4cm] (reduc) at (-4.8, -2.5) {
        \textbf{Reducible}\\
        $\bullet\;\lambda_P = \rho(A) \ge 0 \in \sigma(A)$\\
        $\bullet\;\bx_P \ge \b0,\; \by_P \ge \b0$ (may have $0$s)\\
        $\bullet\;|\lambda_j| \le \lambda_P$
    };
    \node[box, fill=UMbrown!15, font=\bfseries] (irred) at (3.2, -1.5) {Irreducible ($\G(A)$ strongly connected)};
    \node[box, text width=4.6cm, fill=AccentRed!10] (imprim) at (0.4, -4.3) {
        \textbf{Imprimitive (periodic, $d > 1$)}\\
        $\bullet\;\lambda_P > 0$ simple, $\bx_P > \b0, \by_P > \b0$\\
        $\bullet\;|\lambda_j| = \lambda_P$ for $j = 1,\ldots,d$\\
        $\bullet\;|\lambda_j| < \lambda_P$ for $j > d$
    };
    \node[box, text width=4.6cm, fill=LightGreen!25] (prim) at (5.8, -4.3) {
        \textbf{Primitive (aperiodic, $d = 1$)}\\
        $\bullet\;\lambda_P > 0$ simple, $\bx_P > \b0, \by_P > \b0$\\
        $\bullet\;\lambda_P > |\lambda_j|$ for all $j \neq 1$\\
        $\bullet\;(A/\lambda_P)^k$ converges to $\bx_P \by_P^T$
    };
    \draw[->, >=Latex, thick] (nonneg) -- (reduc);
    \draw[->, >=Latex, thick] (nonneg) -- (irred);
    \draw[->, >=Latex, thick] (irred) -- (imprim);
    \draw[->, >=Latex, thick] (irred) -- (prim);
\end{tikzpicture}
\end{center}
\end{frame}

\begin{frame}{Asymptotic behaviour of powers of a primitive matrix}
\begin{importanttheorem}[Asymptotics of primitive matrices]
\label{th:primitive-limit}
Let $A \in \M_n(\IRplus)$ be a primitive matrix with Perron root $\rho(A) > 0$. Let $\bx > \b0$ and $\by > \b0$ be the unique positive right and left eigenvectors satisfying $A\bx = \rho(A)\bx$, $\by^T A = \rho(A)\by^T$, $\11^T\bx = 1$ and $\by^T\bx = 1$. Then
\[
\lim_{k \to \infty} \left(\frac{1}{\rho(A)} A\right)^k = \bx \by^T > 0
\]
\end{importanttheorem}
\vfill
\textbf{Why this holds}: Write $\bz = (\by^T\bz)\bx + \sum_{j=2}^n c_j \bv_j$. Since $|\lambda_j / \rho(A)| < 1$ for all $j \ge 2$,
\[
\left(\frac{1}{\rho(A)} A\right)^k \bz = (\by^T\bz)\bx + \sum_{j=2}^n c_j \left(\frac{\lambda_j}{\rho(A)}\right)^k \bv_j \xrightarrow[k \to \infty]{} (\by^T\bz)\bx = (\bx \by^T)\bz
\]
\end{frame}

\begin{frame}[fragile]{Inspecting the Perron-Frobenius spectrum in R}
Compare the eigenvalues of the periodic matrix $A_1$ ($d = 3$) and primitive matrix $A_2$ ($d = 1$):
\begin{knitrout}
\definecolor{shadecolor}{rgb}{0.969, 0.969, 0.969}\color{fgcolor}\begin{kframe}
\begin{alltt}
\hldef{A1} \hlkwb{<-} \hlkwd{matrix}\hldef{(}\hlkwd{c}\hldef{(}\hlnum{0}\hldef{,} \hlnum{1}\hldef{,} \hlnum{0}\hldef{,}  \hlnum{0}\hldef{,} \hlnum{0}\hldef{,} \hlnum{1}\hldef{,}  \hlnum{1}\hldef{,} \hlnum{0}\hldef{,} \hlnum{0}\hldef{),} \hlkwc{nrow} \hldef{=} \hlnum{3}\hldef{,} \hlkwc{byrow} \hldef{=} \hlnum{TRUE}\hldef{)}
\hldef{A2} \hlkwb{<-} \hlkwd{matrix}\hldef{(}\hlkwd{c}\hldef{(}\hlnum{1}\hldef{,} \hlnum{1}\hldef{,} \hlnum{0}\hldef{,}  \hlnum{0}\hldef{,} \hlnum{0}\hldef{,} \hlnum{1}\hldef{,}  \hlnum{1}\hldef{,} \hlnum{0}\hldef{,} \hlnum{0}\hldef{),} \hlkwc{nrow} \hldef{=} \hlnum{3}\hldef{,} \hlkwc{byrow} \hldef{=} \hlnum{TRUE}\hldef{)}
\hlcom{# Moduli of eigenvalues of A1 (all 3 have modulus 1)}
\hlkwd{round}\hldef{(}\hlkwd{Mod}\hldef{(}\hlkwd{eigen}\hldef{(A1)}\hlopt{$}\hldef{values),} \hlnum{4}\hldef{)}
\end{alltt}
\begin{verbatim}
## [1] 1 1 1
\end{verbatim}
\begin{alltt}
\hlcom{# Eigenvalues and moduli for primitive A2 (unique dominant Perron root)}
\hlkwd{round}\hldef{(}\hlkwd{eigen}\hldef{(A2)}\hlopt{$}\hldef{values,} \hlnum{4}\hldef{)}
\end{alltt}
\begin{verbatim}
## [1]  1.4656+0.0000i -0.2328+0.7926i -0.2328-0.7926i
\end{verbatim}
\begin{alltt}
\hlkwd{round}\hldef{(}\hlkwd{Mod}\hldef{(}\hlkwd{eigen}\hldef{(A2)}\hlopt{$}\hldef{values),} \hlnum{4}\hldef{)}
\end{alltt}
\begin{verbatim}
## [1] 1.4656 0.8260 0.8260
\end{verbatim}
\end{kframe}
\end{knitrout}
\end{frame}


%%%%%%%%%%%%%%%%%%%%
%%%%%%%%%%%%%%%%%%%%
%%%%%%%%%%%%%%%%%%%%
%%%%%%%%%%%%%%%%%%%%
\SSsection{Regular discrete-time Markov chains}

%%%%%%%%%%%%%%%%%%%%
%%%%%%%%%%%%%%%%%%%%
\SSsubsection{Markov chains and column-stochastic matrices}

\begin{frame}{Discrete-time Markov chains}
A \defword{Markov chain} is a discrete-time stochastic process whose evolution depends \textbf{only on the current state} of the system (the \defword{memoryless} or \defword{Markov property}):
\begin{itemize}
\item Consider a system that can be in one of $n$ finite states $S = \{S_1, \ldots, S_n\}$ at discrete times $t = 0, 1, 2, \ldots$
\item The probability of moving to state $S_i$ at time $t+1$ given that the system is in state $S_j$ at time $t$ is
\[
p_{ij} = \IP(S_i(t+1) \mid S_j(t))
\]
\item Note the index order carefully: $p_{ij}$ is the transition probability \textbf{from state $S_j$ (column $j$) to state $S_i$ (row $i$)}
\end{itemize}
\vfill
Markov chains are a cornerstone of data science: they underpin PageRank, Hidden Markov Models (HMMs), Markov Chain Monte Carlo (MCMC) sampling and reinforcement learning
\end{frame}

\begin{frame}{The transition matrix $P$}
\begin{importantdefinition}[Discrete-time Markov chain and transition matrix]
A sequence of random states from $S = \{S_1, \ldots, S_n\}$ is a (time-homogeneous) \defword{discrete-time Markov chain} (DTMC) if for all $t = 0, 1, 2, \ldots$ and all $1 \le i, j \le n$, the transition probability
\[
p_{ij} = \IP(S_i(t+1) \mid S_j(t))
\]
depends only on $i$ and $j$. The matrix $P = [p_{ij}] \in \M_n(\IRplus)$ is called the \defword{transition matrix} of the Markov chain
\end{importantdefinition}
\vfill
If the system is in state $S_j$ at time $t$, it must transition to one of $S_1, \ldots, S_n$ at time $t+1$, so the entries in \textbf{column $j$} of $P$ must sum to $1$:
\[
p_{1j} + p_{2j} + \cdots + p_{nj} = \sum_{i=1}^n p_{ij} = 1 \quad \text{for all } j = 1, \ldots, n
\]
\end{frame}

\begin{frame}{How do we get to state $S_i$ at time $t+1$?}
For $1 \le i \le n$, let $p_i(t) = \IP(S_i(t))$ be the probability that the system is in state $S_i$ at time $t$
\vfill
List all $n$ mutually exclusive ways to arrive in state $S_i$ at time $t+1$:
\begin{enumerate}
\item We were in $S_1$ at time $t$ (probability $p_1(t)$) and transitioned from $S_1$ to $S_i$ (probability $p_{i1}$), contributing $p_{i1} p_1(t)$
\item We were in $S_2$ at time $t$ (probability $p_2(t)$) and transitioned from $S_2$ to $S_i$ (probability $p_{i2}$), contributing $p_{i2} p_2(t)$
\item[$\vdots$]
\item[$n$.] We were in $S_n$ at time $t$ (probability $p_n(t)$) and transitioned from $S_n$ to $S_i$ (probability $p_{in}$), contributing $p_{in} p_n(t)$
\end{enumerate}
\vfill
Summing over all $j = 1, \ldots, n$ gives the law of total probability:
\[
p_i(t+1) = \sum_{j=1}^n p_{ij} p_j(t) = p_{i1}p_1(t) + p_{i2}p_2(t) + \cdots + p_{in}p_n(t)
\]
\end{frame}

\begin{frame}{State distribution update in matrix form}
Collecting the state probabilities into the column \defword{probability vector}
\[
\bp(t) = \begin{pmatrix} p_1(t) \\ \vdots \\ p_n(t) \end{pmatrix} \ge \b0, \qquad \11^T \bp(t) = \sum_{i=1}^n p_i(t) = 1
\]
the update equations for $i = 1, \ldots, n$ take the compact matrix form
\[
\begin{pmatrix}
p_1(t+1) \\ \vdots \\ p_n(t+1)
\end{pmatrix}
=
\begin{pmatrix}
p_{11} & p_{12} & \cdots & p_{1n} \\
\vdots & \vdots & \ddots & \vdots \\
p_{n1} & p_{n2} & \cdots & p_{nn}
\end{pmatrix}
\begin{pmatrix}
p_1(t) \\ \vdots \\ p_n(t)
\end{pmatrix}
\qquad \Longleftrightarrow \qquad
\bp(t+1) = P\bp(t)
\]
\vfill
Iterating from the initial distribution $\bp(0)$ gives $\bp(1) = P\bp(0)$, $\bp(2) = P\bp(1) = P^2\bp(0)$ and for any $t \in \IN$:
\[
\bp(t) = P^t \bp(0)
\]
\end{frame}

\begin{frame}{Transition matrix and transition digraph}
For a column-stochastic transition matrix $P = [p_{ij}]$, recall that $p_{ij} > 0$ represents a transition \textbf{from state $S_j$ to state $S_i$}:
\begin{columns}[c]
\begin{column}{0.46\textwidth}
\centering
\textbf{Column-stochastic matrix $P$}
\[
P = \begin{pmatrix}
0.5 & 0.8 & 0.1 \\
0.5 & 0   & 0.2 \\
0   & 0.2 & 0.7
\end{pmatrix}
\]
\vspace{0.1cm}
\begin{itemize}
\item Column 1 (from $S_1$): $0.5 + 0.5 + 0 = 1$
\item Column 2 (from $S_2$): $0.8 + 0 + 0.2 = 1$
\item Column 3 (from $S_3$): $0.1 + 0.2 + 0.7 = 1$
\end{itemize}
\end{column}
\begin{column}{0.54\textwidth}
\centering
\textbf{Transition digraph}\\[0.2cm]
\begin{tikzpicture}[
    ->,
    >=Latex,
    node distance=2.2cm,
    every path/.style={thick},
    state_node/.style={circle, draw=UMbrown, fill=HighlightGold!25, thick, minimum size=9mm, font=\bfseries}
]
    \node[state_node] (S1) at (0, 2.0) {$S_1$};
    \node[state_node] (S2) at (-2.0, 0) {$S_2$};
    \node[state_node] (S3) at (2.0, 0) {$S_3$};
    \path (S1) edge [loop above] node[above] {$0.5$} (S1);
    \path (S1) edge [bend left=12] node[right] {$0.5$} (S2);
    \path (S2) edge [bend left=12] node[left] {$0.8$} (S1);
    \path (S2) edge [bend left=12] node[above] {$0.2$} (S3);
    \path (S3) edge node[right] {$0.1$} (S1);
    \path (S3) edge [bend left=12] node[below] {$0.2$} (S2);
    \path (S3) edge [loop right] node[right] {$0.7$} (S3);
\end{tikzpicture}
\end{column}
\end{columns}
\vfill
Out-weights from each vertex $S_j$ in the transition digraph correspond to \textbf{column $j$} of $P$ and always sum to $1$
\end{frame}

\begin{frame}{Stochastic matrices}
\begin{importantdefinition}[Column-stochastic, row-stochastic and doubly stochastic matrices]
Let $M = [m_{ij}] \in \M_n(\IRplus)$ be a non-negative matrix and $\11 = (1, \ldots, 1)^T \in \IR^n$
\begin{itemize}
\item $M$ is \defword{column-stochastic} if every column sums to $1$:
\[
\sum_{i=1}^n m_{ij} = 1 \quad \text{for all } j = 1, \ldots, n \qquad \Longleftrightarrow \qquad \11^T M = \11^T
\]
\item $M$ is \defword{row-stochastic} if every row sums to $1$:
\[
\sum_{j=1}^n m_{ij} = 1 \quad \text{for all } i = 1, \ldots, n \qquad \Longleftrightarrow \qquad M\11 = \11
\]
\item $M$ is \defword{doubly stochastic} if it is both column-stochastic and row-stochastic ($\11^T M = \11^T$ and $M\11 = \11$)
\end{itemize}
\end{importantdefinition}
\vfill
Throughout this course, \defword{stochastic matrix} and \defword{transition matrix} mean \textbf{column-stochastic} unless explicitly stated otherwise
\end{frame}

\begin{frame}[red]{Products and powers of column-stochastic matrices}
\begin{theorem}
\label{th:product-stochastic}
\begin{enumerate}
\item If $M, N \in \M_n(\IRplus)$ are column-stochastic matrices, then their product $MN$ is also column-stochastic
\item For any column-stochastic matrix $P \in \M_n(\IRplus)$ and any integer $t \ge 1$, $P^t$ is column-stochastic, and if $\bp(0)$ is a probability vector, then $\bp(t) = P^t\bp(0)$ is a probability vector
\end{enumerate}
\end{theorem}
\vfill
\textbf{Proof}:
\begin{enumerate}
\item Since $M \ge 0$ and $N \ge 0$, their product satisfies $MN \ge 0$. Furthermore, using the associativity of matrix multiplication and $\11^T M = \11^T$, $\11^T N = \11^T$:
\[
\11^T (MN) = (\11^T M) N = \11^T N = \11^T
\]
Thus $MN$ is non-negative and its columns sum to $1$, so $MN$ is column-stochastic
\item By induction on $t \ge 1$, $P^t$ is column-stochastic, and $\bp(t) = P^t \bp(0) \ge \b0$ with $\11^T \bp(t) = (\11^T P^t)\bp(0) = \11^T \bp(0) = 1$ \qed
\end{enumerate}
\end{frame}

\begin{frame}[red]{Spectral radius of a column-stochastic matrix}
\begin{importanttheorem}[Spectrum of a stochastic matrix]
\label{th:spectrum-stochastic}
Let $P \in \M_n(\IRplus)$ be a column-stochastic matrix. Then:
\begin{enumerate}
\item $\lambda_1 = 1$ is an eigenvalue of $P$
\item Every eigenvalue $\lambda \in \sigma(P)$ satisfies $|\lambda| \le 1$, so the spectral radius is $\rho(P) = 1$
\end{enumerate}
\end{importanttheorem}
\vfill
\textbf{Proof}:
\begin{enumerate}
\item Since $P$ is column-stochastic, $\11^T P = \11^T$. Taking the transpose gives $P^T \11 = \11$ with $\11 \neq \b0$, so $\lambda = 1$ is an eigenvalue of $P^T$. Since $\det(P^T - \lambda \II_n) = \det(P - \lambda \II_n)$, $P$ and $P^T$ have the same eigenvalues ($\sigma(P) = \sigma(P^T)$), so $1 \in \sigma(P)$
\item Let $\lambda \in \sigma(P) = \sigma(P^T)$ and let $\bx \neq \b0$ be a right eigenvector of $P^T$: $P^T \bx = \lambda \bx$. Pick an index $k$ maximising $|x_k| = \max_{1 \le j \le n} |x_j| > 0$. Then from row $k$ of $P^T \bx = \lambda \bx$:
\[
|\lambda|\, |x_k| = |\lambda x_k| = \left|\sum_{i=1}^n p_{ik} x_i\right| \le \sum_{i=1}^n p_{ik} |x_i| \le \left(\sum_{i=1}^n p_{ik}\right) |x_k| = 1 \cdot |x_k|
\]
Dividing by $|x_k| > 0$ yields $|\lambda| \le 1$, and since $1 \in \sigma(P)$, we have $\rho(P) = 1$ \qed
\end{enumerate}
\end{frame}


%%%%%%%%%%%%%%%%%%%%
%%%%%%%%%%%%%%%%%%%%
\SSsubsection{The fundamental limit theorem for regular Markov chains}

\begin{frame}{Regular Markov chains}
\begin{importantdefinition}[Regular Markov chain]
A discrete-time Markov chain with column-stochastic transition matrix $P \in \M_n(\IRplus)$ is \defword{regular} if $P$ is \textbf{primitive}, i.e. there exists a positive integer $k \ge 1$ such that
\[
P^k > 0 \quad (\text{every entry of } P^k \text{ is strictly positive})
\]
\end{importantdefinition}
\vfill
\textbf{Interpretation}:
\begin{itemize}
\item The entry $(P^k)_{ij}$ is the $k$-step transition probability of moving from state $S_j$ to state $S_i$ in exactly $k$ steps
\item Thus a Markov chain is regular if and only if there is some time horizon $k \ge 1$ such that from \textbf{any} starting state $S_j$, there is a strictly positive probability of being in \textbf{any} target state $S_i$ after $k$ steps
\item By Theorem~\ref{th:primitive-iff-sc-aperiodic}, a Markov chain is regular $\iff$ its transition digraph is \textbf{strongly connected} (irreducible) and \textbf{aperiodic} ($d = 1$, e.g. at least one state has a self-loop $p_{ii} > 0$)
\end{itemize}
\end{frame}

\begin{frame}{The fundamental limit theorem for regular Markov chains}
Applying Theorem~\ref{th:primitive-limit} with $\rho(P) = 1$ and left Perron vector $\by = \11$ yields:
\begin{importanttheorem}[Fundamental limit theorem for regular Markov chains]
\label{th:regular-mc-limit}
Let $P \in \M_n(\IRplus)$ be the column-stochastic transition matrix of a \textbf{regular} Markov chain. Then:
\begin{enumerate}
\item There exists a \textbf{unique} strictly positive probability vector $\bv = (v_1, \ldots, v_n)^T > \b0$ satisfying $P\bv = \bv$ and $\11^T\bv = \sum_{i=1}^n v_i = 1$, called the \defword{stationary distribution} (or \defword{equilibrium} / \defword{steady-state distribution}) of the Markov chain
\item As $t \to \infty$, the matrix powers $P^t$ converge to the rank-one column-stochastic matrix
\[
\lim_{t \to \infty} P^t = W = \begin{pmatrix} \bv & \bv & \cdots & \bv \end{pmatrix} = \bv \11^T > 0
\]
\item For \textbf{any} initial probability distribution $\bp(0)$, the state vector $\bp(t) = P^t\bp(0)$ converges to $\bv$:
\[
\lim_{t \to \infty} \bp(t) = W\bp(0) = \bv(\11^T \bp(0)) = \bv
\]
\end{enumerate}
\end{importanttheorem}
\end{frame}

\begin{frame}{Finding the stationary distribution $\bv$ in practice}
How do we compute the unique stationary distribution $\bv$ of a regular transition matrix $P \in \M_n(\IRplus)$ on an exam or by hand?
\vfill
\begin{enumerate}
\item \textbf{Write the fixed-point equation $P\bv = \bv$ as a homogeneous linear system}:
\[
(P - \II_n)\bv = \b0 \qquad (\text{so } \bv \in \nullspace(P - \II_n))
\]
\begin{itemize}
\item Since $\lambda_1 = 1$ is a simple eigenvalue of $P$, $\dim\nullspace(P - \II_n) = 1$
\item Because $\11^T(P - \II_n) = \b0^T$, the rows of $P - \II_n$ always sum to zero (one row is redundant)
\end{itemize}
\vfill
\item \textbf{Row-reduce $P - \II_n$} (or solve the first $n-1$ equations) to express $v_1, \ldots, v_n$ in terms of a single free parameter, giving a positive basis vector $\bw \in \nullspace(P - \II_n)$
\vfill
\item \textbf{Normalise to a probability vector} using $\11^T\bv = v_1 + \cdots + v_n = 1$:
\[
\bv = \frac{\bw}{\|\bw\|_1} = \frac{\bw}{w_1 + \cdots + w_n}
\]
\end{enumerate}
\end{frame}


%%%%%%%%%%%%%%%%%%%%
%%%%%%%%%%%%%%%%%%%%
\SSsubsection{Examples and running example: orange cats}

\begin{frame}{Exam-style example 1: general $2 \times 2$ regular Markov chain}
Consider a two-state Markov chain with transition probabilities $\alpha = \IP(S_2 \mid S_1)$ and $\beta = \IP(S_1 \mid S_2)$, where $0 < \alpha, \beta < 1$:
\begin{columns}[c]
\begin{column}{0.45\textwidth}
\centering
\begin{tikzpicture}[
    ->, >=Latex, every path/.style={thick},
    state_node/.style={circle, draw=UMbrown, fill=HighlightGold!25, thick, minimum size=9mm, font=\bfseries}
]
    \node[state_node] (S1) at (0, 0) {$S_1$};
    \node[state_node] (S2) at (2.8, 0) {$S_2$};
    \path (S1) edge [loop left] node[left] {$1-\alpha$} (S1);
    \path (S1) edge [bend left=20] node[above] {$\alpha$} (S2);
    \path (S2) edge [bend left=20] node[below] {$\beta$} (S1);
    \path (S2) edge [loop right] node[right] {$1-\beta$} (S2);
\end{tikzpicture}
\end{column}
\begin{column}{0.55\textwidth}
\[
P = \begin{pmatrix}
1 - \alpha & \beta \\
\alpha & 1 - \beta
\end{pmatrix} > 0
\]
Since $P > 0$, the chain is regular ($P^1 > 0$)
\end{column}
\end{columns}
\vfill
\textbf{Finding the stationary distribution $\bv = (v_1, v_2)^T$}:
\[
(P - \II_2)\bv = \begin{pmatrix} -\alpha & \beta \\ \alpha & -\beta \end{pmatrix} \begin{pmatrix} v_1 \\ v_2 \end{pmatrix} = \begin{pmatrix} 0 \\ 0 \end{pmatrix} \iff -\alpha v_1 + \beta v_2 = 0 \iff v_2 = \frac{\alpha}{\beta} v_1
\]
Imposing $v_1 + v_2 = 1$ gives $v_1\left(1 + \frac{\alpha}{\beta}\right) = 1$, hence:
\[
\bv = \frac{1}{\alpha + \beta} \begin{pmatrix} \beta \\ \alpha \end{pmatrix},
\qquad
\lim_{t \to \infty} P^t = W = \frac{1}{\alpha + \beta} \begin{pmatrix} \beta & \beta \\ \alpha & \alpha \end{pmatrix}
\]
\end{frame}

\begin{frame}{Exam-style example 2: $3 \times 3$ regular chain via row reduction}
\textbf{Problem}: Show that the Markov chain with transition matrix $P$ is regular and find its stationary distribution $\bv$ and $\lim_{t \to \infty} P^t$:
\[
P = \begin{pmatrix}
0.5 & 0.4 & 0 \\
0.5 & 0.2 & 0.5 \\
0   & 0.4 & 0.5
\end{pmatrix}
\]
\vfill
\textbf{Step 1 (Regularity)}:
\begin{itemize}
\item From $P$, we have transitions $S_1 \to S_2 \to S_3 \to S_2 \to S_1$, so all 3 states communicate ($\G(P^T)$ is strongly connected, $P$ is irreducible)
\item Diagonal entries $p_{11} = 0.5 > 0$, $p_{22} = 0.2 > 0$, $p_{33} = 0.5 > 0$ give self-loops, so by Corollary~\ref{cor:self-loop-primitive}, $P$ is \textbf{primitive} (in fact $P^2 > 0$) and the chain is \textbf{regular}
\end{itemize}
\end{frame}

\begin{frame}{Exam-style example 2 (continued): solving $(P - \II_3)\bv = \b0$}
\textbf{Step 2 (Nullspace of $P - \II_3$)}: Form $P - \II_3$ (notice the columns sum to $0$):
{\small
\[
P - \II_3 = \begin{pmatrix}
-0.5 & 0.4 & 0 \\
0.5 & -0.8 & 0.5 \\
0 & 0.4 & -0.5
\end{pmatrix}
\xrightarrow{R_2 + R_1}
\begin{pmatrix}
-0.5 & 0.4 & 0 \\
0 & -0.4 & 0.5 \\
0 & 0.4 & -0.5
\end{pmatrix}
\xrightarrow{R_3 + R_2}
\begin{pmatrix}
-0.5 & 0.4 & 0 \\
0 & -0.4 & 0.5 \\
0 & 0 & 0
\end{pmatrix}
\]
}
From the two pivot rows with $v_3$ as free variable:
\begin{itemize}
\item Row 2: $-0.4 v_2 + 0.5 v_3 = 0 \implies v_2 = \frac{5}{4} v_3$
\item Row 1: $-0.5 v_1 + 0.4 v_2 = 0 \implies v_1 = \frac{4}{5} v_2 = \frac{4}{5} \cdot \frac{5}{4} v_3 = v_3$
\end{itemize}
\textbf{Step 3 (Normalisation)}: Impose $v_1 + v_2 + v_3 = 1$, i.e. $\left(1 + \frac{5}{4} + 1\right) v_3 = \frac{13}{4} v_3 = 1 \implies v_3 = \frac{4}{13}$, giving:
\[
\bv = \begin{pmatrix} 4/13 \\ 5/13 \\ 4/13 \end{pmatrix},
\qquad
\lim_{t \to \infty} P^t = W = \begin{pmatrix}
4/13 & 4/13 & 4/13 \\
5/13 & 5/13 & 5/13 \\
4/13 & 4/13 & 4/13
\end{pmatrix}
\]
\end{frame}

\begin{frame}{Running example: the ``orange'' gene in cats}
A cat's coat colour is determined by many genes; the ``orange'' trait comes from a specific gene called the \textbf{orange locus}
\vfill
It has two \emph{alleles} (versions):
\begin{itemize}
\item \textbf{O} $\implies$ produces \emph{phaeomelanin} (red/orange pigment)
\item \textbf{o} $\implies$ produces \emph{eumelanin} (black/brown pigment)
\end{itemize}
\vfill
This gene is \emph{sex-linked}: it is located on the \textbf{X chromosome}
\begin{itemize}
\item Females have two X chromosomes (\textbf{XX}), receiving one allele from each parent: genotypes $X^O X^O$, $X^o X^o$ or $X^O X^o$
\item Males have one X and one Y chromosome (\textbf{XY}), receiving only one allele (always from the mother): genotypes $X^O Y$ or $X^o Y$
\end{itemize}
\end{frame}

\begin{frame}{Genotype vs phenotype in cats}
\textbf{Males (simple)}:
\begin{itemize}
\item $X^O Y \implies$ \textbf{orange cat}
\item $X^o Y \implies$ \textbf{non-orange cat} (e.g. black)
\end{itemize}
\vfill
\textbf{Females (the special case)}:
\begin{itemize}
\item $X^O X^O \implies$ \textbf{orange cat} ($S_1$)
\item $X^o X^o \implies$ \textbf{non-orange cat} ($S_2$, e.g. black)
\item $X^O X^o \implies$ \textbf{tortoiseshell cat} ($S_3$)
\end{itemize}
\vfill
A ``tortie'' is not a simple blend: because of random X-inactivation in female cells, both alleles ($O$ and $o$) are active in different patches of skin, creating the mottled orange and black pattern (male tortoiseshells are extremely rare genetic anomalies with \textbf{XXY} genotype $X^O X^o Y$)
\end{frame}

\begin{frame}{Punnett square warm-up: crossing two parents}
\begin{columns}[T]
\begin{column}{0.48\textwidth}
\textbf{Cross 1: orange dad ($X^O Y$) + black mum ($X^o X^o$)}
\vfill
\punnett{$X^O$ & $Y$} {
    \textbf{Mother} & $X^o$ & $X^O X^o$ & $X^o Y$ \\ \cmidrule(lr){2-4}
    & $X^o$ & $X^O X^o$ & $X^o Y$
}
\vfill
\begin{itemize}
\item All daughters ($X^O X^o$) are \textbf{tortoiseshell}
\item All sons ($X^o Y$) are \textbf{black}
\end{itemize}
\end{column}
\begin{column}{0.48\textwidth}
\textbf{Cross 2: black dad ($X^o Y$) + tortie mum ($X^O X^o$)}
\vfill
\punnett{$X^o$ & $Y$} {
    \textbf{Mother} & $X^O$ & $X^O X^o$ & $X^O Y$ \\ \cmidrule(lr){2-4}
    & $X^o$ & $X^o X^o$ & $X^o Y$
}
\vfill
\begin{itemize}
\item Daughters: $50\%$ tortie ($X^O X^o$), $50\%$ black ($X^o X^o$)
\item Sons: $50\%$ orange ($X^O Y$), $50\%$ black ($X^o Y$)
\end{itemize}
\end{column}
\end{columns}
\end{frame}

\begin{frame}{Setting up the regular Markov chain for female genotypes}
Track the 3 female genotypes across generations:
\begin{itemize}
\item $S_1$: $X^O X^O$ (orange female)
\item $S_2$: $X^o X^o$ (black female)
\item $S_3$: $X^O X^o$ (tortoiseshell female)
\end{itemize}
\vfill
In each generation, we mate the current female with a male chosen at random from a \textbf{fixed population} consisting of:
\begin{itemize}
\item $50\%$ orange males ($X^O Y$)
\item $50\%$ black males ($X^o Y$)
\end{itemize}
and select one of her daughters at random as the next generation's female
\end{frame}

\begin{frame}{State 1: orange female ($X^O X^O$)}
The mother is $X^O X^O$. We pick a father with $50/50$ probability:
\vfill
\begin{columns}[T]
\begin{column}{.5\textwidth}
\centering
\textbf{Case 1: father is $X^O Y$}\\[0.2cm]
\renewcommand{\arraystretch}{1.5}
\begin{tabular}{c c | c | c |}
\multicolumn{2}{c}{} & \multicolumn{1}{c}{$X^O$} & \multicolumn{1}{c}{$Y$} \\ \cline{3-4}
\multirow{2}{*}{\rotatebox{90}{Mother}} & $X^O$ & \cellcolor{punnettorange}$X^O X^O$ ($S_1$) & \cellcolor{punnettorange}Male \\ \cline{2-4}
& $X^O$ & \cellcolor{punnettorange}$X^O X^O$ ($S_1$) & \cellcolor{punnettorange}Male \\ \cline{2-4}
\end{tabular}
\\[0.2cm]
Daughters: $100\%\; S_1$
\end{column}
\begin{column}{.5\textwidth}
\centering
\textbf{Case 2: father is $X^o Y$}\\[0.2cm]
\renewcommand{\arraystretch}{1.5}
\begin{tabular}{c c | c | c |}
\multicolumn{2}{c}{} & \multicolumn{1}{c}{$X^o$} & \multicolumn{1}{c}{$Y$} \\ \cline{3-4}
\multirow{2}{*}{\rotatebox{90}{Mother}} & $X^O$ & \cellcolor{punnetttortie}$X^O X^o$ ($S_3$) & \cellcolor{punnettorange}Male \\ \cline{2-4}
& $X^O$ & \cellcolor{punnetttortie}$X^O X^o$ ($S_3$) & \cellcolor{punnettorange}Male \\ \cline{2-4}
\end{tabular}
\\[0.2cm]
Daughters: $100\%\; S_3$
\end{column}
\end{columns}
\vfill
\textbf{Transitions from $S_1$ (column 1 of $P$)}:
\begin{itemize}
\item $p_{11} = \IP(S_1 \mid S_1) = 0.5 \times 1.0 = \mathbf{0.5}$
\item $p_{21} = \IP(S_2 \mid S_1) = \mathbf{0}$
\item $p_{31} = \IP(S_3 \mid S_1) = 0.5 \times 1.0 = \mathbf{0.5}$
\end{itemize}
\end{frame}

\begin{frame}{State 2: black female ($X^o X^o$)}
The mother is $X^o X^o$. We pick a father with $50/50$ probability:
\vfill
\begin{columns}[T]
\begin{column}{.5\textwidth}
\centering
\textbf{Case 1: father is $X^O Y$}\\[0.2cm]
\renewcommand{\arraystretch}{1.5}
\begin{tabular}{c c | c | c |}
\multicolumn{2}{c}{} & \multicolumn{1}{c}{$X^O$} & \multicolumn{1}{c}{$Y$} \\ \cline{3-4}
\multirow{2}{*}{\rotatebox{90}{Mother}} & $X^o$ & \cellcolor{punnetttortie}$X^O X^o$ ($S_3$)& \cellcolor{punnettblack}Male \\ \cline{2-4}
& $X^o$ & \cellcolor{punnetttortie}$X^O X^o$ ($S_3$)& \cellcolor{punnettblack}Male \\ \cline{2-4}
\end{tabular}
\\[0.2cm]
Daughters: $100\%\; S_3$
\end{column}
\begin{column}{.5\textwidth}
\centering
\textbf{Case 2: father is $X^o Y$}\\[0.2cm]
\renewcommand{\arraystretch}{1.5}
\begin{tabular}{c c | c | c |}
\multicolumn{2}{c}{} & \multicolumn{1}{c}{$X^o$} & \multicolumn{1}{c}{$Y$} \\ \cline{3-4}
\multirow{2}{*}{\rotatebox{90}{Mother}} & $X^o$ & \cellcolor{punnettblack}$X^o X^o$ ($S_2$)& \cellcolor{punnettblack}Male \\ \cline{2-4}
& $X^o$ & \cellcolor{punnettblack}$X^o X^o$ ($S_2$)& \cellcolor{punnettblack}Male \\ \cline{2-4}
\end{tabular}
\\[0.2cm]
Daughters: $100\%\; S_2$
\end{column}
\end{columns}
\vfill
\textbf{Transitions from $S_2$ (column 2 of $P$)}:
\begin{itemize}
\item $p_{12} = \IP(S_1 \mid S_2) = \mathbf{0}$
\item $p_{22} = \IP(S_2 \mid S_2) = 0.5 \times 1.0 = \mathbf{0.5}$
\item $p_{32} = \IP(S_3 \mid S_2) = 0.5 \times 1.0 = \mathbf{0.5}$
\end{itemize}
\end{frame}

\begin{frame}{State 3: tortoiseshell female ($X^O X^o$)}
The mother is $X^O X^o$. We pick a father with $50/50$ probability:
\vfill
\begin{columns}[T]
\begin{column}{0.48\textwidth}
\centering
\textbf{Case 1: father is $X^O Y$}\\[0.2cm]
\renewcommand{\arraystretch}{1.4}
\begin{tabular}{c c | c | c |}
\multicolumn{2}{c}{} & \multicolumn{1}{c}{$X^O$} & \multicolumn{1}{c}{$Y$} \\ \cline{3-4}
\multirow{2}{*}{\rotatebox{90}{Mother}} & $X^O$ & \cellcolor{punnettorange}$X^O X^O$ ($S_1$) & \cellcolor{punnettorange}Male \\ \cline{2-4}
& $X^o$ & \cellcolor{punnetttortie}$X^O X^o$ ($S_3$) & \cellcolor{punnettblack}Male \\ \cline{2-4}
\end{tabular}
\\[0.2cm]
Daughters: $50\%\; S_1$, $50\%\; S_3$
\end{column}
\begin{column}{0.48\textwidth}
\centering
\textbf{Case 2: father is $X^o Y$}\\[0.2cm]
\renewcommand{\arraystretch}{1.4}
\begin{tabular}{c c | c | c |}
\multicolumn{2}{c}{} & \multicolumn{1}{c}{$X^o$} & \multicolumn{1}{c}{$Y$} \\ \cline{3-4}
\multirow{2}{*}{\rotatebox{90}{Mother}} & $X^O$ & \cellcolor{punnetttortie}$X^O X^o$ ($S_3$) & \cellcolor{punnettorange}Male \\ \cline{2-4}
& $X^o$ & \cellcolor{punnettblack}$X^o X^o$ ($S_2$) & \cellcolor{punnettblack}Male \\ \cline{2-4}
\end{tabular}
\\[0.2cm]
Daughters: $50\%\; S_2$, $50\%\; S_3$
\end{column}
\end{columns}
\vfill
\textbf{Transitions from $S_3$ (column 3 of $P$)}:
\begin{itemize}
\item $p_{13} = \IP(S_1 \mid S_3) = 0.5 \times 0.5 = \mathbf{0.25}$
\item $p_{23} = \IP(S_2 \mid S_3) = 0.5 \times 0.5 = \mathbf{0.25}$
\item $p_{33} = \IP(S_3 \mid S_3) = 0.5 \times 0.5 + 0.5 \times 0.5 = \mathbf{0.5}$
\end{itemize}
\end{frame}

\begin{frame}{Transition matrix and digraph of the orange cat chain}
\begin{columns}[c]
\begin{column}{0.45\textwidth}
\[
P = \begin{pmatrix}
0.5 & 0   & 0.25 \\
0   & 0.5 & 0.25 \\
0.5 & 0.5 & 0.5
\end{pmatrix}
\]
\vspace{0.2cm}
\textbf{Is this Markov chain regular?}
\begin{itemize}
\item \textbf{Strongly connected}: $S_1 \to S_3 \to S_2$ and $S_2 \to S_3 \to S_1$, so all states communicate ($P$ is irreducible)
\item \textbf{Aperiodic}: $p_{11}, p_{22}, p_{33} > 0$ (self-loops), so $P$ is \textbf{primitive} ($P^2 > 0$) and the chain is \textbf{regular}
\end{itemize}
\end{column}
\begin{column}{0.55\textwidth}
\centering
\begin{tikzpicture}[
    ->,
    >=Latex,
    node distance=2.5cm,
    every path/.style={thick}
]
    \node[state, fill=punnettorange] (S1) at (0, 2.1) {$X^OX^O$};
    \node[state, fill=punnettblack, text=white] (S2) at (-2.1, 0) {$X^oX^o$};
    \node[state, fill=punnetttortie, text=white] (S3) at (2.1, 0) {$X^OX^o$};
    \path (S1) edge [loop left] node[left] {$0.5$} (S1);
    \path (S1) edge [bend left=12] node[above, sloped] {$0.5$} (S3);
    \path (S2) edge [loop left] node[left] {$0.5$} (S2);
    \path (S2) edge [bend left=12] node[above] {$0.5$} (S3);
    \path (S3) edge [loop right] node[right] {$0.5$} (S3);
    \path (S3) edge [bend left=12] node[sloped, below] {$0.25$} (S1);
    \path (S3) edge [bend left=12] node[below] {$0.25$} (S2);
\end{tikzpicture}
\end{column}
\end{columns}
\end{frame}

\begin{frame}{Stationary distribution of the orange cat chain}
We solve $(P - \II_3)\bv = \b0$ with $v_1 + v_2 + v_3 = 1$:
\[
(P - \II_3)\bv = \begin{pmatrix}
-0.5 & 0    & 0.25 \\
0    & -0.5 & 0.25 \\
0.5  & 0.5  & -0.5
\end{pmatrix}
\begin{pmatrix} v_1 \\ v_2 \\ v_3 \end{pmatrix}
= \begin{pmatrix} 0 \\ 0 \\ 0 \end{pmatrix}
\]
From the first two equations:
\begin{itemize}
\item $-0.5 v_1 + 0.25 v_3 = 0 \implies v_1 = 0.5 v_3$
\item $-0.5 v_2 + 0.25 v_3 = 0 \implies v_2 = 0.5 v_3$
\end{itemize}
Substituting into the normalisation condition $v_1 + v_2 + v_3 = 1$:
\[
0.5 v_3 + 0.5 v_3 + v_3 = 2 v_3 = 1 \implies v_3 = 0.5, \quad v_1 = 0.25, \quad v_2 = 0.25
\]
\vfill
Thus the unique stationary distribution is $\bv = (0.25, 0.25, 0.5)^T$: in the long run, regardless of the initial female's genotype, $25\%$ of females are orange, $25\%$ are black and $50\%$ are tortoiseshell
\end{frame}

\begin{frame}[fragile]{Computing matrix powers and stationary distributions in R}
Let us verify the convergence of $P^t$ and the eigenpair $P\bv = \bv$ in \code{R}:
\begin{knitrout}
\definecolor{shadecolor}{rgb}{0.969, 0.969, 0.969}\color{fgcolor}\begin{kframe}
\begin{alltt}
\hldef{P_reg} \hlkwb{<-} \hlkwd{matrix}\hldef{(}\hlkwd{c}\hldef{(}\hlnum{0.5}\hldef{,} \hlnum{0}\hldef{,}   \hlnum{0.25}\hldef{,}
                  \hlnum{0}\hldef{,}   \hlnum{0.5}\hldef{,} \hlnum{0.25}\hldef{,}
                  \hlnum{0.5}\hldef{,} \hlnum{0.5}\hldef{,} \hlnum{0.5}\hldef{),}
                \hlkwc{nrow} \hldef{=} \hlnum{3}\hldef{,} \hlkwc{byrow} \hldef{=} \hlnum{TRUE}\hldef{)}
\hlcom{# Dominant right eigenvector corresponding to eigenvalue 1}
\hldef{ev} \hlkwb{<-} \hlkwd{eigen}\hldef{(P_reg)}
\hldef{v_stat} \hlkwb{<-} \hlkwd{Re}\hldef{(ev}\hlopt{$}\hldef{vectors[,} \hlnum{1}\hldef{])}
\hldef{v_stat} \hlkwb{<-} \hldef{v_stat} \hlopt{/} \hlkwd{sum}\hldef{(v_stat)}
\hldef{v_stat}
\end{alltt}
\begin{verbatim}
## [1] 0.25 0.25 0.50
\end{verbatim}
\begin{alltt}
\hlcom{# Compare with P^20 computed by repeated multiplication}
\hldef{P_20} \hlkwb{<-} \hlkwd{Reduce}\hldef{(`%*%`,} \hlkwd{replicate}\hldef{(}\hlnum{20}\hldef{, P_reg,} \hlkwc{simplify} \hldef{=} \hlnum{FALSE}\hldef{))}
\hlkwd{round}\hldef{(P_20,} \hlnum{6}\hldef{)}
\end{alltt}
\begin{verbatim}
##      [,1] [,2] [,3]
## [1,] 0.25 0.25 0.25
## [2,] 0.25 0.25 0.25
## [3,] 0.50 0.50 0.50
\end{verbatim}
\end{kframe}
\end{knitrout}
\end{frame}

\begin{frame}[fragile]{Convergence of $\bp(t) = P^t \bp(0)$ from an initial state}
\begin{knitrout}
\definecolor{shadecolor}{rgb}{0.969, 0.969, 0.969}\color{fgcolor}\begin{kframe}
\begin{alltt}
\hldef{t_max} \hlkwb{<-} \hlnum{12}
\hldef{p_hist} \hlkwb{<-} \hlkwd{matrix}\hldef{(}\hlnum{0}\hldef{,} \hlkwc{nrow} \hldef{=} \hlnum{3}\hldef{,} \hlkwc{ncol} \hldef{= t_max} \hlopt{+} \hlnum{1}\hldef{)}
\hldef{p_hist[,} \hlnum{1}\hldef{]} \hlkwb{<-} \hlkwd{c}\hldef{(}\hlnum{1}\hldef{,} \hlnum{0}\hldef{,} \hlnum{0}\hldef{)} \hlcom{# Start with 100% orange female (S_1)}
\hlkwa{for} \hldef{(t} \hlkwa{in} \hlnum{1}\hlopt{:}\hldef{t_max) \{}
  \hldef{p_hist[, t} \hlopt{+} \hlnum{1}\hldef{]} \hlkwb{<-} \hldef{P_reg} \hlopt{%*%} \hldef{p_hist[, t]}
\hldef{\}}
\hlkwd{matplot}\hldef{(}\hlnum{0}\hlopt{:}\hldef{t_max,} \hlkwd{t}\hldef{(p_hist),} \hlkwc{type} \hldef{=} \hlsng{"b"}\hldef{,} \hlkwc{pch} \hldef{=} \hlnum{19}\hldef{,} \hlkwc{lty} \hldef{=} \hlnum{1}\hldef{,} \hlkwc{lwd} \hldef{=} \hlnum{2}\hldef{,}
        \hlkwc{col} \hldef{=} \hlkwd{c}\hldef{(}\hlsng{"darkorange"}\hldef{,} \hlsng{"black"}\hldef{,} \hlsng{"brown"}\hldef{),}
        \hlkwc{xlab} \hldef{=} \hlsng{"Generation t"}\hldef{,} \hlkwc{ylab} \hldef{=} \hlsng{"Probability p_i(t)"}\hldef{,}
        \hlkwc{ylim} \hldef{=} \hlkwd{c}\hldef{(}\hlnum{0}\hldef{,} \hlnum{1}\hldef{),} \hlkwc{main} \hldef{=} \hlsng{"Convergence to stationary distribution v = (0.25, 0.25, 0.5)"}\hldef{)}
\hlkwd{legend}\hldef{(}\hlsng{"topright"}\hldef{,} \hlkwc{legend} \hldef{=} \hlkwd{c}\hldef{(}\hlsng{"S1 (orange)"}\hldef{,} \hlsng{"S2 (black)"}\hldef{,} \hlsng{"S3 (tortie)"}\hldef{),}
       \hlkwc{col} \hldef{=} \hlkwd{c}\hldef{(}\hlsng{"darkorange"}\hldef{,} \hlsng{"black"}\hldef{,} \hlsng{"brown"}\hldef{),} \hlkwc{lty} \hldef{=} \hlnum{1}\hldef{,} \hlkwc{pch} \hldef{=} \hlnum{19}\hldef{,} \hlkwc{lwd} \hldef{=} \hlnum{2}\hldef{)}
\end{alltt}
\end{kframe}
\end{knitrout}
\end{frame}

\maxFrameImage{FIGS/U08-u08-convergence-plot-1.pdf}

\begin{frame}[fragile]{Simulating a single realisation of the regular Markov chain}
In a column-stochastic chain, when in state $j$ at step $t$, the next state at step $t+1$ is drawn with probabilities given by \textbf{column $j$} of $P$:
\begin{knitrout}
\definecolor{shadecolor}{rgb}{0.969, 0.969, 0.969}\color{fgcolor}\begin{kframe}
\begin{alltt}
\hlkwd{set.seed}\hldef{(}\hlnum{2740}\hldef{)}
\hldef{n_steps} \hlkwb{<-} \hlnum{40}
\hldef{states} \hlkwb{<-} \hlkwd{numeric}\hldef{(n_steps} \hlopt{+} \hlnum{1}\hldef{)}
\hldef{states[}\hlnum{1}\hldef{]} \hlkwb{<-} \hlnum{1} \hlcom{# Start in state S_1 (orange female)}
\hlkwa{for} \hldef{(t} \hlkwa{in} \hlnum{1}\hlopt{:}\hldef{n_steps) \{}
  \hldef{states[t} \hlopt{+} \hlnum{1}\hldef{]} \hlkwb{<-} \hlkwd{sample}\hldef{(}\hlnum{1}\hlopt{:}\hlnum{3}\hldef{,} \hlkwc{size} \hldef{=} \hlnum{1}\hldef{,} \hlkwc{prob} \hldef{= P_reg[, states[t]])}
\hldef{\}}
\hlkwd{plot}\hldef{(}\hlnum{0}\hlopt{:}\hldef{n_steps, states,} \hlkwc{type} \hldef{=} \hlsng{"b"}\hldef{,} \hlkwc{pch} \hldef{=} \hlnum{19}\hldef{,} \hlkwc{lwd} \hldef{=} \hlnum{1.8}\hldef{,} \hlkwc{col} \hldef{=} \hlsng{"darkblue"}\hldef{,}
     \hlkwc{yaxt} \hldef{=} \hlsng{"n"}\hldef{,} \hlkwc{xlab} \hldef{=} \hlsng{"Generation t"}\hldef{,} \hlkwc{ylab} \hldef{=} \hlsng{"State"}\hldef{,}
     \hlkwc{main} \hldef{=} \hlsng{"Single realisation of the orange cat Markov chain (40 generations)"}\hldef{)}
\hlkwd{axis}\hldef{(}\hlnum{2}\hldef{,} \hlkwc{at} \hldef{=} \hlnum{1}\hlopt{:}\hlnum{3}\hldef{,} \hlkwc{labels} \hldef{=} \hlkwd{c}\hldef{(}\hlsng{"S1 (orange)"}\hldef{,} \hlsng{"S2 (black)"}\hldef{,} \hlsng{"S3 (tortie)"}\hldef{))}
\end{alltt}
\end{kframe}
\end{knitrout}
\end{frame}

\maxFrameImage{FIGS/U08-u08-dtmc-realisation-1.pdf}

\begin{frame}[fragile]{Empirical state frequencies vs stationary distribution $\bv$}
By the ergodic theorem for regular Markov chains, the long-run fraction of time spent in each state along a single trajectory converges to $\bv = (0.25, 0.25, 0.5)^T$:
\begin{knitrout}
\definecolor{shadecolor}{rgb}{0.969, 0.969, 0.969}\color{fgcolor}\begin{kframe}
\begin{alltt}
\hlkwd{set.seed}\hldef{(}\hlnum{2740}\hldef{)}
\hldef{n_long} \hlkwb{<-} \hlnum{10000}
\hldef{traj} \hlkwb{<-} \hlkwd{numeric}\hldef{(n_long)}
\hldef{traj[}\hlnum{1}\hldef{]} \hlkwb{<-} \hlnum{1}
\hlkwa{for} \hldef{(t} \hlkwa{in} \hlnum{1}\hlopt{:}\hldef{(n_long} \hlopt{-} \hlnum{1}\hldef{)) \{}
  \hldef{traj[t} \hlopt{+} \hlnum{1}\hldef{]} \hlkwb{<-} \hlkwd{sample}\hldef{(}\hlnum{1}\hlopt{:}\hlnum{3}\hldef{,} \hlkwc{size} \hldef{=} \hlnum{1}\hldef{,} \hlkwc{prob} \hldef{= P_reg[, traj[t]])}
\hldef{\}}
\hlcom{# Empirical frequencies over 10,000 generations vs exact v}
\hlkwd{rbind}\hldef{(}\hlkwc{empirical} \hldef{=} \hlkwd{prop.table}\hldef{(}\hlkwd{table}\hldef{(}\hlkwd{factor}\hldef{(traj,} \hlkwc{levels} \hldef{=} \hlnum{1}\hlopt{:}\hlnum{3}\hldef{))),}
      \hlkwc{stationary} \hldef{= v_stat)}
\end{alltt}
\begin{verbatim}
##                 1      2      3
## empirical  0.2414 0.2505 0.5081
## stationary 0.2500 0.2500 0.5000
\end{verbatim}
\end{kframe}
\end{knitrout}
\end{frame}




\end{document}
